% =====================================================================
% CHAPTER 3 — SYSTEM DESIGN
% Graduation project: a policy-based secure governance layer for
% autonomous OS-level agents, built as a hard fork of OpenClaw.
% Kinan Radaideh, Mohammad Al-Masri, Malek Tluli. PSUT, Spring 2025-2026.
%
% STATUS: OUTLINE (2026-09-20). Every heading carries a one-line
%   placeholder saying what the finished section will argue. Replace
%   each placeholder with the real prose; the \sectionstub macro is
%   defined below so every unwritten section is easy to find, and
%   removing the last \sectionstub means the chapter is done.
%
% ---------------------------------------------------------------------
% WHY EVERY HEADING HAS TEXT UNDER IT
% ---------------------------------------------------------------------
% LaTeX will not break a page immediately after a sectioning heading —
% it sets a penalty so a heading is never left orphaned at the foot of a
% page. When headings follow one another with nothing but comments
% between them, there is no legal breakpoint anywhere in the run, the
% whole block becomes unbreakable, and it overflows the bottom of the
% page with an "Overfull \vbox" warning instead of continuing overleaf.
% Measured on 2026-09-20: forty bare headings compile to 2 pages with an
% overfull vbox; the same forty with one line of text each compile to 4
% pages with none. Comments set no type, so they do not count.
% Keep at least one line of real text under every heading.
%
% ---------------------------------------------------------------------
% HOW TO USE THIS FILE
% ---------------------------------------------------------------------
% Paste the body below into main.tex after Chapter 2, or keep it here
% and pull it in with % =====================================================================
% CHAPTER 3 — SYSTEM DESIGN
% Graduation project: a policy-based secure governance layer for
% autonomous OS-level agents, built as a hard fork of OpenClaw.
% Kinan Radaideh, Mohammad Al-Masri, Malek Tluli. PSUT, Spring 2025-2026.
%
% STATUS: OUTLINE (2026-09-20). Every heading carries a one-line
%   placeholder saying what the finished section will argue. Replace
%   each placeholder with the real prose; the \sectionstub macro is
%   defined below so every unwritten section is easy to find, and
%   removing the last \sectionstub means the chapter is done.
%
% ---------------------------------------------------------------------
% WHY EVERY HEADING HAS TEXT UNDER IT
% ---------------------------------------------------------------------
% LaTeX will not break a page immediately after a sectioning heading —
% it sets a penalty so a heading is never left orphaned at the foot of a
% page. When headings follow one another with nothing but comments
% between them, there is no legal breakpoint anywhere in the run, the
% whole block becomes unbreakable, and it overflows the bottom of the
% page with an "Overfull \vbox" warning instead of continuing overleaf.
% Measured on 2026-09-20: forty bare headings compile to 2 pages with an
% overfull vbox; the same forty with one line of text each compile to 4
% pages with none. Comments set no type, so they do not count.
% Keep at least one line of real text under every heading.
%
% ---------------------------------------------------------------------
% HOW TO USE THIS FILE
% ---------------------------------------------------------------------
% Paste the body below into main.tex after Chapter 2, or keep it here
% and pull it in with % =====================================================================
% CHAPTER 3 — SYSTEM DESIGN
% Graduation project: a policy-based secure governance layer for
% autonomous OS-level agents, built as a hard fork of OpenClaw.
% Kinan Radaideh, Mohammad Al-Masri, Malek Tluli. PSUT, Spring 2025-2026.
%
% STATUS: OUTLINE (2026-09-20). Every heading carries a one-line
%   placeholder saying what the finished section will argue. Replace
%   each placeholder with the real prose; the \sectionstub macro is
%   defined below so every unwritten section is easy to find, and
%   removing the last \sectionstub means the chapter is done.
%
% ---------------------------------------------------------------------
% WHY EVERY HEADING HAS TEXT UNDER IT
% ---------------------------------------------------------------------
% LaTeX will not break a page immediately after a sectioning heading —
% it sets a penalty so a heading is never left orphaned at the foot of a
% page. When headings follow one another with nothing but comments
% between them, there is no legal breakpoint anywhere in the run, the
% whole block becomes unbreakable, and it overflows the bottom of the
% page with an "Overfull \vbox" warning instead of continuing overleaf.
% Measured on 2026-09-20: forty bare headings compile to 2 pages with an
% overfull vbox; the same forty with one line of text each compile to 4
% pages with none. Comments set no type, so they do not count.
% Keep at least one line of real text under every heading.
%
% ---------------------------------------------------------------------
% HOW TO USE THIS FILE
% ---------------------------------------------------------------------
% Paste the body below into main.tex after Chapter 2, or keep it here
% and pull it in with % =====================================================================
% CHAPTER 3 — SYSTEM DESIGN
% Graduation project: a policy-based secure governance layer for
% autonomous OS-level agents, built as a hard fork of OpenClaw.
% Kinan Radaideh, Mohammad Al-Masri, Malek Tluli. PSUT, Spring 2025-2026.
%
% STATUS: OUTLINE (2026-09-20). Every heading carries a one-line
%   placeholder saying what the finished section will argue. Replace
%   each placeholder with the real prose; the \sectionstub macro is
%   defined below so every unwritten section is easy to find, and
%   removing the last \sectionstub means the chapter is done.
%
% ---------------------------------------------------------------------
% WHY EVERY HEADING HAS TEXT UNDER IT
% ---------------------------------------------------------------------
% LaTeX will not break a page immediately after a sectioning heading —
% it sets a penalty so a heading is never left orphaned at the foot of a
% page. When headings follow one another with nothing but comments
% between them, there is no legal breakpoint anywhere in the run, the
% whole block becomes unbreakable, and it overflows the bottom of the
% page with an "Overfull \vbox" warning instead of continuing overleaf.
% Measured on 2026-09-20: forty bare headings compile to 2 pages with an
% overfull vbox; the same forty with one line of text each compile to 4
% pages with none. Comments set no type, so they do not count.
% Keep at least one line of real text under every heading.
%
% ---------------------------------------------------------------------
% HOW TO USE THIS FILE
% ---------------------------------------------------------------------
% Paste the body below into main.tex after Chapter 2, or keep it here
% and pull it in with \input{chapter3}. The \sectionstub definition has
% to go in the preamble; it is repeated at the top of the body so the
% file compiles on its own, and can be deleted once the chapter is done.
%
% The chapter is titled "System Design" because that is what §1.8
% (Document Organization) promises the reader.
%
% ---------------------------------------------------------------------
% HOUSE STYLE FOR THIS CHAPTER
% ---------------------------------------------------------------------
% 1. HEADINGS ARE NOUN PHRASES, two to four words, title case. Not
%    sentences and not questions, following the two model reports in
%    Documentation/GradProj/, whose Chapter 3 headings read "Packet
%    Handling", "Voting Process", "System Architecture", "Smart
%    Contract", "Donor Setup", "System Security". The argument goes in
%    the prose, not the heading.
% 2. AMERICAN SPELLING, because Chapters 1 and 2 already use it:
%    "prioritize", "utilizes", "Data Sanitization", "organizations",
%    "behavior". So: canonicalization, sanitization, organization,
%    authorization, behavior. The repository's own notes are British
%    English; do not carry that into the report.
% 3. The section spine (requirements, analysis of requirements, analysis
%    of constraints, different design approaches, then Developed Design
%    last before the summary) is the spine both model reports use, and
%    the one CHAPTER3-MATERIAL.md is keyed to.
%
% ---------------------------------------------------------------------
% WHERE THE BUILT DESIGN DIFFERS FROM CHAPTERS 1 AND 2
% ---------------------------------------------------------------------
% It is normal and expected that the final implementation differs from
% the preliminary design. What matters is that Chapter 3 states each
% difference openly AND explains why the change was made. The register
% of them, with the justification for each, is
% docs-notes/report/DIVERGENCES.md. Sections that owe an explanation
% carry a % DIVERGENCE comment pointing at it. Add to that file when a
% new one is found, and add the reasoning, not just the fact.
%
% ---------------------------------------------------------------------
% FIGURES: how they are referenced, and where to find them
% ---------------------------------------------------------------------
% The drawings are NOT in this file. They live in docs-notes/FIGURES.md,
% which holds each figure three ways (prose, Mermaid, TikZ) with the
% reasoning behind it. That file is the source of truth.
%
% Prose cites a figure by the label FIGURES.md already gives it, so the
% two files cannot drift:
%
%     ...as shown in Figure~\ref{fig:gov-architecture}.
%
% Every Chapter 3 and 4 figure label is prefixed "fig:gov-". The prefix
% is deliberate: Chapter 1 already uses \label{fig:architecture} for the
% preliminary design figure, and without the prefix our F1 would have
% silently collided with it. Keep the prefix for anything new. Tables
% take "tab:gov-" for the same reason.
%
% TO SEE a figure:  open docs-notes/figures/report-figures.pdf at the
%                   page in the table below.
% TO GET its TikZ:  copy it from docs-notes/FIGURES.md under
%                   "## F1: ..." -> "### TikZ form", or lift the whole
%                   figure environment out of the generated file
%                   docs-notes/figures/report-figures.tex.
% TO REBUILD both:  node docs-notes/figures/build-figures.mjs
%                   then run pdflatex on report-figures.tex.
%
% All fourteen compile clean and have been checked by eye (2026-09-20).
%
%   Prints as | FIGURES.md | Cite with                    | PDF page
%   ----------+------------+------------------------------+---------
%   Fig 3.1   | F1         | \ref{fig:gov-architecture}   |  1
%   Fig 3.2   | F2         | \ref{fig:gov-rbac}           |  2
%   Fig 3.3   | F3         | \ref{fig:gov-decision}       |  3
%   Fig 3.4   | F5         | \ref{fig:gov-pathnorm}       |  4
%   Fig 3.5   | F6         | \ref{fig:gov-promptpath}     |  5
%   Fig 3.6   | F8         | \ref{fig:gov-twopaths}       |  6
%   Fig 3.7   | F11        | \ref{fig:gov-toctou}         |  7
%   Fig 3.8   | F19        | \ref{fig:gov-tenant}         |  8
%   Fig 3.9   | F21        | \ref{fig:gov-codexgates}     |  9
%   Fig 3.10  | F22        | \ref{fig:gov-foldergrant}    | 10
%   Fig 3.11  | F23        | \ref{fig:gov-alwaysallow}    | 11
%   Fig 3.12  | F24        | \ref{fig:gov-taskslot}       | 12
%   Fig 4.1   | F14        | \ref{fig:gov-coverage}       | 13  (Chapter 4)
%   Fig 4.2   | F17        | \ref{fig:gov-defectage}      | 14  (Chapter 4)
%
% main.tex already has tikz, the \usetikzlibrary line, pgfplots with
% compat=1.18, and graphicx (via style-PSUT-BSc-Eng). It still needs the
% shared \tikzset style block from FIGURES.md -> "One shared style
% block", pasted after \pgfplotsset{compat=1.18}. Without it the figures
% compile but lose every box, arrow and label style.
% =====================================================================

% Marks an unwritten section. Put this in the main.tex preamble; delete
% it, and every use of it, once the chapter is written.
\providecommand{\sectionstub}[1]{\emph{[Outline: #1]}\par}

\chapter{System Design}
\label{ch:design}

The preceding chapter established that OS-level Agents introduce a kind of
vulnerability that is cognitive rather than deterministic, and that what an
attacker often manipulates is the reasoning of the agent itself. Any control
that lives inside the model, or that relies on the model's own judgment to
enforce it, therefore inherits the weakness it is meant to correct. The design
in this chapter, briefly described in previous sections, starts with the
assumption that the agent is untrusted, and a deterministic policy check is
placed on the path that every tool call by an agent takes. This required that
the governance layer is built as a hard fork compiled into the OpenClaw core
rather than as a plugin installed alongside it. This commits the project to an
entire upstream codebase, including native client applications that this work
does not modify.

On that foundation the layer adds five capabilities to the host: a default-deny
policy gate on the tool-call path, a tamper-evident audit ledger, an emergency
kill-switch that measures the time to a confirmed stop, four tiers of named
human accounts with individual sessions and levels of authority, and a
management interface through which each of these is operated. The chapter moves
from requirement to finished system. Section~\ref{sec:gov-requirements} briefly
restates the nine design requirements of Section~1.3.
Section~\ref{sec:gov-reqanalysis} then analyzes them, describing the status the
system has reached in relation to each one.
Section~\ref{sec:gov-constraints} does the same for the design constraints of
Section~1.4, and Section~\ref{sec:gov-standards} describes how the engineering
standards of Section~1.5 are applied. Section~\ref{sec:gov-approaches} presents
the alternatives that were seriously considered and the reason each was
rejected, and Section~\ref{sec:gov-developed} describes the system as built, one
design area at a time, ending with what the system does not defend against; in
other words, accepted limitations. Where the developed design departs from the
preliminary design of Chapter~1, the departure is stated where it arises and the
reason is given with it.

% =====================================================================
\section{Design Requirements}
\label{sec:gov-requirements}

As stated in Chapter 1, the design shall achieve the following requirements:

\begin{enumerate}
    \item The system shall be implemented using the Node.js runtime (version 18 or higher) and developed primarily in TypeScript with static type checking enabled.
    \item The system shall provide a secure web-based dashboard that enables administrators to configure customized privilege policies, monitor active autonomous agent sessions, and manage system operations using role-based access control.
    \item The system shall enforce a default-deny, policy-based security model that restricts autonomous agent access to operating system resources, including file system paths, process execution, and network communication, based on administrator-defined capabilities.
    \item The system shall support customized, fine-grained privileges for autonomous agents, including path-level file access, command allowlisting, network allowlisting, and time-limited permissions.
    \item The system shall continuously monitor autonomous agent activities and shall record 100\% of agent actions, policy decisions, and administrative approvals in an auditable logging system.
    \item The system shall implement tamper-evident audit logging mechanisms to ensure the integrity and traceability of all recorded agent and administrative actions.
    \item The system shall provide real-time administrative control capabilities, including the ability to suspend or terminate an autonomous agent session within a maximum response time of one second.
    \item The system shall prevent sensitive data (such as secrets or credentials) from being written in plaintext to log files.
    \item The system shall be deployable on a Linux-based operating system using open-source software components only.
\end{enumerate}


% =====================================================================
\section{Analysis of Design Requirements}
\label{sec:gov-reqanalysis}

Section~\ref{sec:gov-requirements} states the requirements as specified. This
section examines what each one demands of the design when it is read strictly,
and records the status the implementation has reached against it. The
distinction is material, because each requirement is a single sentence of
specification and several of them constrain the design considerably more
tightly than their length suggests. Two turn on the interpretation of a single
phrase, and a third was satisfied only after its initial interpretation was
discarded and the implementation replaced. Table~\ref{tab:gov-reqstatus}
summarizes the outcome and indicates where each requirement is discussed. All
nine requirements are met. The ninth, deployment on a Linux host, was the last
to close, and the evidence that closed it is given in
Section~\ref{sec:gov-remaining-reqs}.

\begin{table}[h]
\centering
\caption{Design requirement status}
\label{tab:gov-reqstatus}
\renewcommand{\arraystretch}{1.3}
\small
\begin{tabular}{@{}c p{5.7cm} l l@{}}
\toprule
\textbf{\#} & \textbf{Requirement} & \textbf{Status} & \textbf{Discussed in} \\
\midrule
1 & Node.js and TypeScript, with static type checking & Met & \S\ref{sec:gov-remaining-reqs} \\
2 & Secure web dashboard with role-based access control & Met & \S\ref{sec:gov-dashboard} \\
3 & Default-deny model over paths, execution and network & Met & \S\ref{sec:gov-enforcement} \\
4 & Fine-grained, time-limited privileges & Met & \S\ref{sec:gov-gate} \\
5 & Record 100\% of actions, decisions and approvals & Met & \S\ref{sec:gov-coverage} \\
6 & Tamper-evident audit logging & Met & \S\ref{sec:gov-ledger} \\
7 & Suspend or terminate a session within one second & Met & \S\ref{sec:gov-latency} \\
8 & No plaintext secrets in log files & Met & \S\ref{sec:gov-credentials} \\
9 & Deployable on Linux, open-source components only & Met & \S\ref{sec:gov-remaining-reqs} \\
\bottomrule
\end{tabular}
\end{table}

Four requirements are treated individually below, because in each case a strict
reading either altered the design or altered what could honestly be claimed of
it. Requirement 3 turned on the position of the check rather than on the
decision it reaches. Requirement 5 turned on the scope of \emph{every action},
and its first implementation was replaced. Requirement 7 turned on the quantity
being measured, the dispatch of a stop signal or the stop itself. Requirement 8
was satisfied by reusing an existing component rather than implementing a second
one, which is a design decision and is presented as such. The remaining five
requirements are addressed together in Section~\ref{sec:gov-remaining-reqs}.

\subsection{Enforcement Point}
\label{sec:gov-enforcement}
\sectionstub{Requirement 3. Why a default-deny gate has to sit on the
host's own tool-call path rather than beside it, and what that costs in
coverage.}

\subsection{Audit Coverage}
\label{sec:gov-coverage}
\sectionstub{Requirement 5. What recording one hundred percent of agent
actions was first read as meaning, why that reading was abandoned, and
the distinct decision that lets an unevaluated call be recorded as such
rather than folded in with the allowed ones.}

\subsection{Termination Latency}
\label{sec:gov-latency}
\sectionstub{Requirement 7. The difference between the stop signal being
sent and the stop being confirmed, and which of the two the requirement
is read as claiming.}

\subsection{Credential Protection}
\label{sec:gov-credentials}
\sectionstub{Requirement 8. Met by redacting at the ledger boundary with
the host's own redaction rather than by reimplementing it.}

\subsection{Remaining Requirements}
\label{sec:gov-remaining-reqs}
\sectionstub{Requirements 1, 2, 4, 6 and 9 in brief: what each asked,
and what satisfying it implied for the design.}

% =====================================================================
\section{Analysis of Design Constraints}
\label{sec:gov-constraints}

\sectionstub{The three categories of Section 1.4, and what each one
ruled out of the design.}

% MATERIAL: CHAPTER3-MATERIAL.md "→ 3.3 Analysis of Design Constraints".

\subsection{Economic Constraints}
\sectionstub{The 350 JOD ceiling and the open-source-only stack, met in
part by adding no new dependencies at all, which is a claim any reader
can check.}

\subsection{Manufacturability and Sustainability}
\sectionstub{What the layer adds to resident memory and to startup time
on an 8\,GB server, and the startup budget that was set and defended
because of this constraint.}

\subsection{Ethical and Professional Responsibility}
\sectionstub{The design can subtract capability from an agent and never
add it, and why that asymmetry was made structural rather than left to
care.}

\subsection{Accepted Limitations}
\sectionstub{Constraints the design accepted rather than solved:
English-only operator text, and development on Windows against a Linux
deployment target.}

% =====================================================================
\section{Discussion of Engineering Standards}
\label{sec:gov-standards}

\sectionstub{How the design applies the two standards named in Section
1.5, rather than restating what those standards are.}

% NOTE: the one section with no material behind it in the repository.
%   Write it from §1.5's wording plus the design as built.

\subsection{ISO/IEC 27001:2022}
\sectionstub{Policy-driven access control mapped onto the default-deny
gate and its rule language; systematic governance of rights onto the
four tiers and ownership; audit logging onto the hash-chained ledger;
and risk treatment onto the baseline policy and the core denials that
cannot be switched off at runtime.}

\subsection{OWASP Secure Coding Practices}
\sectionstub{Session management and authentication mapped onto the two
gates on every governance request; error handling onto refusals that
explain themselves without leaking; data in transit onto keeping the
dashboard off the public internet behind an SSH tunnel; and input
handling onto rule validation and path canonicalization.}

% =====================================================================
\section{Different Design Approaches}
\label{sec:gov-approaches}

\sectionstub{The alternatives that were genuinely weighed, each stated
as the options and then the choice, with the reason the others lost.}

% MATERIAL: CHAPTER3-MATERIAL.md "→ 3.4" holds two of these in full; the
%   rest are recoverable from the §3.5 sections named under each heading.

\subsection{Integration Strategy}
\sectionstub{A fork compiled into the host's core, or a plugin beside
it. Why a plugin cannot make a gate unavoidable.}
% MATERIAL: §3.5.2b.

\subsection{Gate Placement}
\sectionstub{Where in the tool-call path the check belongs, and what
each candidate position would and would not have seen.}
% MATERIAL: §3.5.1, §3.5.25.

\subsection{Host Interception}
\sectionstub{How the host is told that the gate must be consulted: four
options, one chosen.}
% MATERIAL: "→ 3.4" §3.4.y, §3.5.15.

\subsection{Path Representation}
\sectionstub{Which form a file path takes when a rule is matched against
it: three options, one chosen.}
% MATERIAL: "→ 3.4" §3.4.x, §3.5.8.

\subsection{Log Integrity}
\sectionstub{A hash chain weighed against the alternatives, and what
each would have cost in complexity and in trust.}
% MATERIAL: §3.5.3, §3.5.9.

\subsection{Operator Identity}
\sectionstub{One shared Gateway credential, or a second and independent
set of named accounts with their own sessions.}
% MATERIAL: §3.5.6, §3.5.19.

\subsection{Management Surfaces}
\sectionstub{Three surfaces were built and one was removed, and why
subtraction counts as a design decision rather than a retreat.}
% MATERIAL: §3.5.77.

\subsection{Tenancy Model}
\sectionstub{How many organizations one installation holds. One, decided
deliberately, with the structure for several deliberately kept.}
% MATERIAL: §3.5.89.

% =====================================================================
\section{Developed Design}
\label{sec:gov-developed}

\sectionstub{The system as built, taken one design area at a time. Each
part is introduced by what it is for, then how it works, then the one
decision inside it that a reader would otherwise stop and question.}

\subsection{System Architecture}
\label{sec:gov-arch}
\sectionstub{The layer inside the Gateway process; the two doors into
it, one for the operator and one for agent activity; and the files it
owns on disk. Figure~\ref{fig:gov-architecture}.}
% FIGURE F1 -> \ref{fig:gov-architecture} (prints as Figure 3.1)
% MATERIAL: §3.5.1, §3.5.2, §3.5.2b.

\subsection{Policy Engine}
\label{sec:gov-gate}
\sectionstub{The component that turns a rule set and an attempted action
into a decision.}

\subsubsection{Rule Model}
\sectionstub{Tiers, effects, access directions, resource kinds and
expiry: the vocabulary a rule is written in.}
% MATERIAL: §3.5.4, §3.5.12, §3.5.20; docs-notes/PERMISSION-SPEC.md.

\subsubsection{Evaluation Order}
\sectionstub{Why denials are read before allowances, and what that
guarantees. Figure~\ref{fig:gov-decision}.}
% FIGURE F3 -> \ref{fig:gov-decision} (prints as Figure 3.3)
% MATERIAL: §3.5.5, §3.5.25.
% DIVERGENCE from Chapter 1 (report/DIVERGENCES.md §1.3): §1.6 says
%   conflicts are resolved by preferring the earlier rule. That is true
%   of the conflict DETECTOR but not of evaluation, where a denial beats
%   an allowance whatever the order. Complete the picture here.

\subsubsection{Path Canonicalization}
\sectionstub{Turning what an agent wrote into the path it resolves to on
disk, before any rule is matched against it.
Figure~\ref{fig:gov-pathnorm}.}
% FIGURE F5 -> \ref{fig:gov-pathnorm} (prints as Figure 3.4)
% MATERIAL: §3.5.8, §3.5.29.

\subsubsection{Baseline Policy}
\sectionstub{What an installation ships with: core denials that cannot
be edited at runtime, and the allowances that make an agent useful
before anyone has written a rule.}
% MATERIAL: docs-notes/BASELINE-RULES.md.

\subsubsection{Folder Grants}
\sectionstub{Granting a directory with exceptions carved out of it, and
why the denial is written first. Figure~\ref{fig:gov-foldergrant}.}
% FIGURE F22 -> \ref{fig:gov-foldergrant} (prints as Figure 3.10)
% MATERIAL: §3.5.66.

\subsection{Audit Ledger}
\label{sec:gov-ledger}
\sectionstub{The append-only record of what was attempted, what was
decided and who decided it.}

\subsubsection{Entry Structure}
\sectionstub{What a single entry holds, including the raw model intent
that produced the action.}
% MATERIAL: §3.5.3, §3.5.59.

\subsubsection{Hash Chaining and Verification}
\sectionstub{How each entry is bound to the one before it, and the
independent verifier that detects any later edit.}
% MATERIAL: §3.5.3; scripts/verify-ledger.mjs.
% DIVERGENCE from Chapter 2 (report/DIVERGENCES.md §1.1): Chapter 2
%   describes a Merkle tree; a keyed HMAC-SHA256 chain was built. State
%   the difference AND why: the tree's advantage is proving one entry
%   without reading the log, which no party here needs, and the keying
%   is a strengthening over Chapter 2's plain SHA-256.

\subsubsection{Administrative Logging}
\sectionstub{Why administrative actions are recorded in the same chain
as agent actions, and what each such entry must name.}
% MATERIAL: §3.5.9, §3.5.63.

\subsubsection{Data Sanitization}
\sectionstub{What an audit trail is allowed not to see, and where
redaction is applied so that it cannot be bypassed.}
% MATERIAL: §3.5.28, §3.5.60.
% DIVERGENCE from Chapter 2 (report/DIVERGENCES.md §1.2): Chapter 2
%   promises regex and entropy-based masking; what was built reuses the
%   host's own redactor at the ledger boundary. State the difference AND
%   why: one tested redactor rather than two to maintain, and placement
%   at the single write boundary beats cleverness applied per caller.

\subsection{Access Control Model}
\label{sec:gov-rbac}
\sectionstub{Who the operators are, what each of them may do, and how
the system decides. Figure~\ref{fig:gov-rbac}.}
% FIGURE F2 -> \ref{fig:gov-rbac} (prints as Figure 3.2)

\subsubsection{Role Hierarchy}
\sectionstub{The four tiers, what each inherits from the one below, and
the capability that deliberately does not sit at the top.}
% MATERIAL: §3.5.4, §3.5.24; docs-notes/ROLE-MODEL.md.

\subsubsection{Two-Gate Authentication}
\sectionstub{The Gateway credential, and then the governance account
session and its role, on every request.}
% MATERIAL: §3.5.6, §3.5.19.

\subsubsection{Ownership and Assignment}
\sectionstub{What managing an agent resolves to in code: the two scope
questions each tier answers independently.}
% MATERIAL: docs-notes/ROLE-MODEL.md, §3.5.22, §3.5.30.

\subsection{Prompt Execution Path}
\label{sec:gov-promptpath}
\sectionstub{A User prompting an agent they hold: the checks made, and
the record written before the run rather than after it.
Figure~\ref{fig:gov-promptpath}.}
% FIGURE F6 -> \ref{fig:gov-promptpath} (prints as Figure 3.5)
% MATERIAL: §3.5.11, §3.5.17.

\subsection{Session Control}
\label{sec:gov-killswitch}
\sectionstub{The two ways an operator intervenes in a running agent, and
what each one is able to reach.}

\subsubsection{Agent Lockdown}
\sectionstub{What a lockdown covers, including work the locked agent had
already started, and why it outranks any allowing rule.}
% MATERIAL: §3.5.38, §3.5.16.

\subsubsection{Kill Switch}
\sectionstub{Stopping a run, and why the design reports a confirmed stop
rather than a signal that was sent.}
% MATERIAL: §3.5.10, §3.5.21.
% DIVERGENCE from Chapter 1 (report/DIVERGENCES.md §1.4): §1.6 promises
%   a sub-second kill switch. The signal is milliseconds; the CONFIRMED
%   stop measured 2.2 s and 2.8 s on the laptop. Give both numbers and
%   say the design measures something stricter than Chapter 1 did.

\subsection{Human-in-the-Loop Approvals}
\label{sec:gov-approvals}
\sectionstub{What happens when the policy neither allows nor denies an
action and a person has to decide.}

\subsubsection{Escalation Routing}
\sectionstub{Who is entitled to answer a request, where it appears, and
what happens when nobody answers in time.}
% MATERIAL: §3.5.83.
% DIVERGENCE from Chapter 1 (report/DIVERGENCES.md §1.5): §1.6 puts the
%   timeout with Root and adds a per-user toggle. What was built sets
%   escalation, posture and timeout per AGENT, by an Administrator, and
%   has no per-user toggle.

\subsubsection{Persistent Approvals}
\sectionstub{What the ``always allow'' answer promises, what it does in
practice, and the rule request it files for an administrator.
Figure~\ref{fig:gov-alwaysallow}.}
% FIGURE F23 -> \ref{fig:gov-alwaysallow} (prints as Figure 3.11)
% MATERIAL: §3.5.81, §3.5.94.

\subsection{Runtime Management}
\label{sec:gov-runs}
\sectionstub{How running work is represented, bounded and governed,
including when the agent executes outside this process.}

\subsubsection{Task and Slot Model}
\sectionstub{The row an operator sees and the concurrency slot behind
it, and why the two are released at different moments.
Figure~\ref{fig:gov-taskslot}.}
% FIGURE F24 -> \ref{fig:gov-taskslot} (prints as Figure 3.12)
% MATERIAL: §3.5.82.

\subsubsection{Secondary Runtime Governance}
\sectionstub{The in-process path and the native harness, and the
two-layer permission that governs the second.
Figures~\ref{fig:gov-twopaths} and~\ref{fig:gov-codexgates}.}
% FIGURE F8  -> \ref{fig:gov-twopaths}   (prints as Figure 3.6)
% FIGURE F21 -> \ref{fig:gov-codexgates} (prints as Figure 3.9)
% MATERIAL: §3.5.15, §3.5.62.

\subsection{Tenancy and Agent Registry}
\label{sec:gov-tenant}
\sectionstub{One organization per installation, the storage that belongs
to it, and the registry that gives the layer a noun for an agent.
Figure~\ref{fig:gov-tenant}.}
% FIGURE F19 -> \ref{fig:gov-tenant} (prints as Figure 3.8)
% MATERIAL: §3.5.31, §3.5.33, §3.5.47, §3.5.89.

\subsection{Agent Lifecycle}
\label{sec:gov-lifecycle}
\sectionstub{Registration as the precondition for being allowed
anything; ownership and re-owning; and the two deletions, with what each
one leaves behind on the host.}
% MATERIAL: §3.5.33, §3.5.90, §3.5.91, §3.5.94.

\subsection{Management Interface}
\label{sec:gov-dashboard}
\sectionstub{The HTTP control plane and the dashboard behind it, what
its sections are for, and the rule that a control must never promise
what the gate cannot keep.}
% MATERIAL: §3.5.13, §3.5.37, §3.5.46, §3.5.80.

\subsection{System Security}
\label{sec:gov-security}
\sectionstub{The threats this design answers, and the three it does not:
the check-then-open window, a denied file still reachable on the native
harness, and the fact that the layer governs what an agent does rather
than why it does it. Figure~\ref{fig:gov-toctou}.}
% FIGURE F11 -> \ref{fig:gov-toctou} (prints as Figure 3.7)
% MATERIAL: §3.5.7, §3.5.61; mg/HANDOFF.md §7.
% CAUTION: do not write that the layer prevents prompt injection. It
%   contains the damage when persuasion succeeds. Chapter 2 sets this
%   up, so this section should answer Chapter 2 directly.

% =====================================================================
\section{Summary}
\label{sec:gov-designsummary}

\sectionstub{What the chapter established, the one requirement left
partially met, and the handover to Chapter 4, which tests the design
against Section 1.3 as its benchmark.}
. The \sectionstub definition has
% to go in the preamble; it is repeated at the top of the body so the
% file compiles on its own, and can be deleted once the chapter is done.
%
% The chapter is titled "System Design" because that is what §1.8
% (Document Organization) promises the reader.
%
% ---------------------------------------------------------------------
% HOUSE STYLE FOR THIS CHAPTER
% ---------------------------------------------------------------------
% 1. HEADINGS ARE NOUN PHRASES, two to four words, title case. Not
%    sentences and not questions, following the two model reports in
%    Documentation/GradProj/, whose Chapter 3 headings read "Packet
%    Handling", "Voting Process", "System Architecture", "Smart
%    Contract", "Donor Setup", "System Security". The argument goes in
%    the prose, not the heading.
% 2. AMERICAN SPELLING, because Chapters 1 and 2 already use it:
%    "prioritize", "utilizes", "Data Sanitization", "organizations",
%    "behavior". So: canonicalization, sanitization, organization,
%    authorization, behavior. The repository's own notes are British
%    English; do not carry that into the report.
% 3. The section spine (requirements, analysis of requirements, analysis
%    of constraints, different design approaches, then Developed Design
%    last before the summary) is the spine both model reports use, and
%    the one CHAPTER3-MATERIAL.md is keyed to.
%
% ---------------------------------------------------------------------
% WHERE THE BUILT DESIGN DIFFERS FROM CHAPTERS 1 AND 2
% ---------------------------------------------------------------------
% It is normal and expected that the final implementation differs from
% the preliminary design. What matters is that Chapter 3 states each
% difference openly AND explains why the change was made. The register
% of them, with the justification for each, is
% docs-notes/report/DIVERGENCES.md. Sections that owe an explanation
% carry a % DIVERGENCE comment pointing at it. Add to that file when a
% new one is found, and add the reasoning, not just the fact.
%
% ---------------------------------------------------------------------
% FIGURES: how they are referenced, and where to find them
% ---------------------------------------------------------------------
% The drawings are NOT in this file. They live in docs-notes/FIGURES.md,
% which holds each figure three ways (prose, Mermaid, TikZ) with the
% reasoning behind it. That file is the source of truth.
%
% Prose cites a figure by the label FIGURES.md already gives it, so the
% two files cannot drift:
%
%     ...as shown in Figure~\ref{fig:gov-architecture}.
%
% Every Chapter 3 and 4 figure label is prefixed "fig:gov-". The prefix
% is deliberate: Chapter 1 already uses \label{fig:architecture} for the
% preliminary design figure, and without the prefix our F1 would have
% silently collided with it. Keep the prefix for anything new. Tables
% take "tab:gov-" for the same reason.
%
% TO SEE a figure:  open docs-notes/figures/report-figures.pdf at the
%                   page in the table below.
% TO GET its TikZ:  copy it from docs-notes/FIGURES.md under
%                   "## F1: ..." -> "### TikZ form", or lift the whole
%                   figure environment out of the generated file
%                   docs-notes/figures/report-figures.tex.
% TO REBUILD both:  node docs-notes/figures/build-figures.mjs
%                   then run pdflatex on report-figures.tex.
%
% All fourteen compile clean and have been checked by eye (2026-09-20).
%
%   Prints as | FIGURES.md | Cite with                    | PDF page
%   ----------+------------+------------------------------+---------
%   Fig 3.1   | F1         | \ref{fig:gov-architecture}   |  1
%   Fig 3.2   | F2         | \ref{fig:gov-rbac}           |  2
%   Fig 3.3   | F3         | \ref{fig:gov-decision}       |  3
%   Fig 3.4   | F5         | \ref{fig:gov-pathnorm}       |  4
%   Fig 3.5   | F6         | \ref{fig:gov-promptpath}     |  5
%   Fig 3.6   | F8         | \ref{fig:gov-twopaths}       |  6
%   Fig 3.7   | F11        | \ref{fig:gov-toctou}         |  7
%   Fig 3.8   | F19        | \ref{fig:gov-tenant}         |  8
%   Fig 3.9   | F21        | \ref{fig:gov-codexgates}     |  9
%   Fig 3.10  | F22        | \ref{fig:gov-foldergrant}    | 10
%   Fig 3.11  | F23        | \ref{fig:gov-alwaysallow}    | 11
%   Fig 3.12  | F24        | \ref{fig:gov-taskslot}       | 12
%   Fig 4.1   | F14        | \ref{fig:gov-coverage}       | 13  (Chapter 4)
%   Fig 4.2   | F17        | \ref{fig:gov-defectage}      | 14  (Chapter 4)
%
% main.tex already has tikz, the \usetikzlibrary line, pgfplots with
% compat=1.18, and graphicx (via style-PSUT-BSc-Eng). It still needs the
% shared \tikzset style block from FIGURES.md -> "One shared style
% block", pasted after \pgfplotsset{compat=1.18}. Without it the figures
% compile but lose every box, arrow and label style.
% =====================================================================

% Marks an unwritten section. Put this in the main.tex preamble; delete
% it, and every use of it, once the chapter is written.
\providecommand{\sectionstub}[1]{\emph{[Outline: #1]}\par}

\chapter{System Design}
\label{ch:design}

The preceding chapter established that OS-level Agents introduce a kind of
vulnerability that is cognitive rather than deterministic, and that what an
attacker often manipulates is the reasoning of the agent itself. Any control
that lives inside the model, or that relies on the model's own judgment to
enforce it, therefore inherits the weakness it is meant to correct. The design
in this chapter, briefly described in previous sections, starts with the
assumption that the agent is untrusted, and a deterministic policy check is
placed on the path that every tool call by an agent takes. This required that
the governance layer is built as a hard fork compiled into the OpenClaw core
rather than as a plugin installed alongside it. This commits the project to an
entire upstream codebase, including native client applications that this work
does not modify.

On that foundation the layer adds five capabilities to the host: a default-deny
policy gate on the tool-call path, a tamper-evident audit ledger, an emergency
kill-switch that measures the time to a confirmed stop, four tiers of named
human accounts with individual sessions and levels of authority, and a
management interface through which each of these is operated. The chapter moves
from requirement to finished system. Section~\ref{sec:gov-requirements} briefly
restates the nine design requirements of Section~1.3.
Section~\ref{sec:gov-reqanalysis} then analyzes them, describing the status the
system has reached in relation to each one.
Section~\ref{sec:gov-constraints} does the same for the design constraints of
Section~1.4, and Section~\ref{sec:gov-standards} describes how the engineering
standards of Section~1.5 are applied. Section~\ref{sec:gov-approaches} presents
the alternatives that were seriously considered and the reason each was
rejected, and Section~\ref{sec:gov-developed} describes the system as built, one
design area at a time, ending with what the system does not defend against; in
other words, accepted limitations. Where the developed design departs from the
preliminary design of Chapter~1, the departure is stated where it arises and the
reason is given with it.

% =====================================================================
\section{Design Requirements}
\label{sec:gov-requirements}

As stated in Chapter 1, the design shall achieve the following requirements:

\begin{enumerate}
    \item The system shall be implemented using the Node.js runtime (version 18 or higher) and developed primarily in TypeScript with static type checking enabled.
    \item The system shall provide a secure web-based dashboard that enables administrators to configure customized privilege policies, monitor active autonomous agent sessions, and manage system operations using role-based access control.
    \item The system shall enforce a default-deny, policy-based security model that restricts autonomous agent access to operating system resources, including file system paths, process execution, and network communication, based on administrator-defined capabilities.
    \item The system shall support customized, fine-grained privileges for autonomous agents, including path-level file access, command allowlisting, network allowlisting, and time-limited permissions.
    \item The system shall continuously monitor autonomous agent activities and shall record 100\% of agent actions, policy decisions, and administrative approvals in an auditable logging system.
    \item The system shall implement tamper-evident audit logging mechanisms to ensure the integrity and traceability of all recorded agent and administrative actions.
    \item The system shall provide real-time administrative control capabilities, including the ability to suspend or terminate an autonomous agent session within a maximum response time of one second.
    \item The system shall prevent sensitive data (such as secrets or credentials) from being written in plaintext to log files.
    \item The system shall be deployable on a Linux-based operating system using open-source software components only.
\end{enumerate}


% =====================================================================
\section{Analysis of Design Requirements}
\label{sec:gov-reqanalysis}

Section~\ref{sec:gov-requirements} states the requirements as specified. This
section examines what each one demands of the design when it is read strictly,
and records the status the implementation has reached against it. The
distinction is material, because each requirement is a single sentence of
specification and several of them constrain the design considerably more
tightly than their length suggests. Two turn on the interpretation of a single
phrase, and a third was satisfied only after its initial interpretation was
discarded and the implementation replaced. Table~\ref{tab:gov-reqstatus}
summarizes the outcome and indicates where each requirement is discussed. All
nine requirements are met. The ninth, deployment on a Linux host, was the last
to close, and the evidence that closed it is given in
Section~\ref{sec:gov-remaining-reqs}.

\begin{table}[h]
\centering
\caption{Design requirement status}
\label{tab:gov-reqstatus}
\renewcommand{\arraystretch}{1.3}
\small
\begin{tabular}{@{}c p{5.7cm} l l@{}}
\toprule
\textbf{\#} & \textbf{Requirement} & \textbf{Status} & \textbf{Discussed in} \\
\midrule
1 & Node.js and TypeScript, with static type checking & Met & \S\ref{sec:gov-remaining-reqs} \\
2 & Secure web dashboard with role-based access control & Met & \S\ref{sec:gov-dashboard} \\
3 & Default-deny model over paths, execution and network & Met & \S\ref{sec:gov-enforcement} \\
4 & Fine-grained, time-limited privileges & Met & \S\ref{sec:gov-gate} \\
5 & Record 100\% of actions, decisions and approvals & Met & \S\ref{sec:gov-coverage} \\
6 & Tamper-evident audit logging & Met & \S\ref{sec:gov-ledger} \\
7 & Suspend or terminate a session within one second & Met & \S\ref{sec:gov-latency} \\
8 & No plaintext secrets in log files & Met & \S\ref{sec:gov-credentials} \\
9 & Deployable on Linux, open-source components only & Met & \S\ref{sec:gov-remaining-reqs} \\
\bottomrule
\end{tabular}
\end{table}

Four requirements are treated individually below, because in each case a strict
reading either altered the design or altered what could honestly be claimed of
it. Requirement 3 turned on the position of the check rather than on the
decision it reaches. Requirement 5 turned on the scope of \emph{every action},
and its first implementation was replaced. Requirement 7 turned on the quantity
being measured, the dispatch of a stop signal or the stop itself. Requirement 8
was satisfied by reusing an existing component rather than implementing a second
one, which is a design decision and is presented as such. The remaining five
requirements are addressed together in Section~\ref{sec:gov-remaining-reqs}.

\subsection{Enforcement Point}
\label{sec:gov-enforcement}
\sectionstub{Requirement 3. Why a default-deny gate has to sit on the
host's own tool-call path rather than beside it, and what that costs in
coverage.}

\subsection{Audit Coverage}
\label{sec:gov-coverage}
\sectionstub{Requirement 5. What recording one hundred percent of agent
actions was first read as meaning, why that reading was abandoned, and
the distinct decision that lets an unevaluated call be recorded as such
rather than folded in with the allowed ones.}

\subsection{Termination Latency}
\label{sec:gov-latency}
\sectionstub{Requirement 7. The difference between the stop signal being
sent and the stop being confirmed, and which of the two the requirement
is read as claiming.}

\subsection{Credential Protection}
\label{sec:gov-credentials}
\sectionstub{Requirement 8. Met by redacting at the ledger boundary with
the host's own redaction rather than by reimplementing it.}

\subsection{Remaining Requirements}
\label{sec:gov-remaining-reqs}
\sectionstub{Requirements 1, 2, 4, 6 and 9 in brief: what each asked,
and what satisfying it implied for the design.}

% =====================================================================
\section{Analysis of Design Constraints}
\label{sec:gov-constraints}

\sectionstub{The three categories of Section 1.4, and what each one
ruled out of the design.}

% MATERIAL: CHAPTER3-MATERIAL.md "→ 3.3 Analysis of Design Constraints".

\subsection{Economic Constraints}
\sectionstub{The 350 JOD ceiling and the open-source-only stack, met in
part by adding no new dependencies at all, which is a claim any reader
can check.}

\subsection{Manufacturability and Sustainability}
\sectionstub{What the layer adds to resident memory and to startup time
on an 8\,GB server, and the startup budget that was set and defended
because of this constraint.}

\subsection{Ethical and Professional Responsibility}
\sectionstub{The design can subtract capability from an agent and never
add it, and why that asymmetry was made structural rather than left to
care.}

\subsection{Accepted Limitations}
\sectionstub{Constraints the design accepted rather than solved:
English-only operator text, and development on Windows against a Linux
deployment target.}

% =====================================================================
\section{Discussion of Engineering Standards}
\label{sec:gov-standards}

\sectionstub{How the design applies the two standards named in Section
1.5, rather than restating what those standards are.}

% NOTE: the one section with no material behind it in the repository.
%   Write it from §1.5's wording plus the design as built.

\subsection{ISO/IEC 27001:2022}
\sectionstub{Policy-driven access control mapped onto the default-deny
gate and its rule language; systematic governance of rights onto the
four tiers and ownership; audit logging onto the hash-chained ledger;
and risk treatment onto the baseline policy and the core denials that
cannot be switched off at runtime.}

\subsection{OWASP Secure Coding Practices}
\sectionstub{Session management and authentication mapped onto the two
gates on every governance request; error handling onto refusals that
explain themselves without leaking; data in transit onto keeping the
dashboard off the public internet behind an SSH tunnel; and input
handling onto rule validation and path canonicalization.}

% =====================================================================
\section{Different Design Approaches}
\label{sec:gov-approaches}

\sectionstub{The alternatives that were genuinely weighed, each stated
as the options and then the choice, with the reason the others lost.}

% MATERIAL: CHAPTER3-MATERIAL.md "→ 3.4" holds two of these in full; the
%   rest are recoverable from the §3.5 sections named under each heading.

\subsection{Integration Strategy}
\sectionstub{A fork compiled into the host's core, or a plugin beside
it. Why a plugin cannot make a gate unavoidable.}
% MATERIAL: §3.5.2b.

\subsection{Gate Placement}
\sectionstub{Where in the tool-call path the check belongs, and what
each candidate position would and would not have seen.}
% MATERIAL: §3.5.1, §3.5.25.

\subsection{Host Interception}
\sectionstub{How the host is told that the gate must be consulted: four
options, one chosen.}
% MATERIAL: "→ 3.4" §3.4.y, §3.5.15.

\subsection{Path Representation}
\sectionstub{Which form a file path takes when a rule is matched against
it: three options, one chosen.}
% MATERIAL: "→ 3.4" §3.4.x, §3.5.8.

\subsection{Log Integrity}
\sectionstub{A hash chain weighed against the alternatives, and what
each would have cost in complexity and in trust.}
% MATERIAL: §3.5.3, §3.5.9.

\subsection{Operator Identity}
\sectionstub{One shared Gateway credential, or a second and independent
set of named accounts with their own sessions.}
% MATERIAL: §3.5.6, §3.5.19.

\subsection{Management Surfaces}
\sectionstub{Three surfaces were built and one was removed, and why
subtraction counts as a design decision rather than a retreat.}
% MATERIAL: §3.5.77.

\subsection{Tenancy Model}
\sectionstub{How many organizations one installation holds. One, decided
deliberately, with the structure for several deliberately kept.}
% MATERIAL: §3.5.89.

% =====================================================================
\section{Developed Design}
\label{sec:gov-developed}

\sectionstub{The system as built, taken one design area at a time. Each
part is introduced by what it is for, then how it works, then the one
decision inside it that a reader would otherwise stop and question.}

\subsection{System Architecture}
\label{sec:gov-arch}
\sectionstub{The layer inside the Gateway process; the two doors into
it, one for the operator and one for agent activity; and the files it
owns on disk. Figure~\ref{fig:gov-architecture}.}
% FIGURE F1 -> \ref{fig:gov-architecture} (prints as Figure 3.1)
% MATERIAL: §3.5.1, §3.5.2, §3.5.2b.

\subsection{Policy Engine}
\label{sec:gov-gate}
\sectionstub{The component that turns a rule set and an attempted action
into a decision.}

\subsubsection{Rule Model}
\sectionstub{Tiers, effects, access directions, resource kinds and
expiry: the vocabulary a rule is written in.}
% MATERIAL: §3.5.4, §3.5.12, §3.5.20; docs-notes/PERMISSION-SPEC.md.

\subsubsection{Evaluation Order}
\sectionstub{Why denials are read before allowances, and what that
guarantees. Figure~\ref{fig:gov-decision}.}
% FIGURE F3 -> \ref{fig:gov-decision} (prints as Figure 3.3)
% MATERIAL: §3.5.5, §3.5.25.
% DIVERGENCE from Chapter 1 (report/DIVERGENCES.md §1.3): §1.6 says
%   conflicts are resolved by preferring the earlier rule. That is true
%   of the conflict DETECTOR but not of evaluation, where a denial beats
%   an allowance whatever the order. Complete the picture here.

\subsubsection{Path Canonicalization}
\sectionstub{Turning what an agent wrote into the path it resolves to on
disk, before any rule is matched against it.
Figure~\ref{fig:gov-pathnorm}.}
% FIGURE F5 -> \ref{fig:gov-pathnorm} (prints as Figure 3.4)
% MATERIAL: §3.5.8, §3.5.29.

\subsubsection{Baseline Policy}
\sectionstub{What an installation ships with: core denials that cannot
be edited at runtime, and the allowances that make an agent useful
before anyone has written a rule.}
% MATERIAL: docs-notes/BASELINE-RULES.md.

\subsubsection{Folder Grants}
\sectionstub{Granting a directory with exceptions carved out of it, and
why the denial is written first. Figure~\ref{fig:gov-foldergrant}.}
% FIGURE F22 -> \ref{fig:gov-foldergrant} (prints as Figure 3.10)
% MATERIAL: §3.5.66.

\subsection{Audit Ledger}
\label{sec:gov-ledger}
\sectionstub{The append-only record of what was attempted, what was
decided and who decided it.}

\subsubsection{Entry Structure}
\sectionstub{What a single entry holds, including the raw model intent
that produced the action.}
% MATERIAL: §3.5.3, §3.5.59.

\subsubsection{Hash Chaining and Verification}
\sectionstub{How each entry is bound to the one before it, and the
independent verifier that detects any later edit.}
% MATERIAL: §3.5.3; scripts/verify-ledger.mjs.
% DIVERGENCE from Chapter 2 (report/DIVERGENCES.md §1.1): Chapter 2
%   describes a Merkle tree; a keyed HMAC-SHA256 chain was built. State
%   the difference AND why: the tree's advantage is proving one entry
%   without reading the log, which no party here needs, and the keying
%   is a strengthening over Chapter 2's plain SHA-256.

\subsubsection{Administrative Logging}
\sectionstub{Why administrative actions are recorded in the same chain
as agent actions, and what each such entry must name.}
% MATERIAL: §3.5.9, §3.5.63.

\subsubsection{Data Sanitization}
\sectionstub{What an audit trail is allowed not to see, and where
redaction is applied so that it cannot be bypassed.}
% MATERIAL: §3.5.28, §3.5.60.
% DIVERGENCE from Chapter 2 (report/DIVERGENCES.md §1.2): Chapter 2
%   promises regex and entropy-based masking; what was built reuses the
%   host's own redactor at the ledger boundary. State the difference AND
%   why: one tested redactor rather than two to maintain, and placement
%   at the single write boundary beats cleverness applied per caller.

\subsection{Access Control Model}
\label{sec:gov-rbac}
\sectionstub{Who the operators are, what each of them may do, and how
the system decides. Figure~\ref{fig:gov-rbac}.}
% FIGURE F2 -> \ref{fig:gov-rbac} (prints as Figure 3.2)

\subsubsection{Role Hierarchy}
\sectionstub{The four tiers, what each inherits from the one below, and
the capability that deliberately does not sit at the top.}
% MATERIAL: §3.5.4, §3.5.24; docs-notes/ROLE-MODEL.md.

\subsubsection{Two-Gate Authentication}
\sectionstub{The Gateway credential, and then the governance account
session and its role, on every request.}
% MATERIAL: §3.5.6, §3.5.19.

\subsubsection{Ownership and Assignment}
\sectionstub{What managing an agent resolves to in code: the two scope
questions each tier answers independently.}
% MATERIAL: docs-notes/ROLE-MODEL.md, §3.5.22, §3.5.30.

\subsection{Prompt Execution Path}
\label{sec:gov-promptpath}
\sectionstub{A User prompting an agent they hold: the checks made, and
the record written before the run rather than after it.
Figure~\ref{fig:gov-promptpath}.}
% FIGURE F6 -> \ref{fig:gov-promptpath} (prints as Figure 3.5)
% MATERIAL: §3.5.11, §3.5.17.

\subsection{Session Control}
\label{sec:gov-killswitch}
\sectionstub{The two ways an operator intervenes in a running agent, and
what each one is able to reach.}

\subsubsection{Agent Lockdown}
\sectionstub{What a lockdown covers, including work the locked agent had
already started, and why it outranks any allowing rule.}
% MATERIAL: §3.5.38, §3.5.16.

\subsubsection{Kill Switch}
\sectionstub{Stopping a run, and why the design reports a confirmed stop
rather than a signal that was sent.}
% MATERIAL: §3.5.10, §3.5.21.
% DIVERGENCE from Chapter 1 (report/DIVERGENCES.md §1.4): §1.6 promises
%   a sub-second kill switch. The signal is milliseconds; the CONFIRMED
%   stop measured 2.2 s and 2.8 s on the laptop. Give both numbers and
%   say the design measures something stricter than Chapter 1 did.

\subsection{Human-in-the-Loop Approvals}
\label{sec:gov-approvals}
\sectionstub{What happens when the policy neither allows nor denies an
action and a person has to decide.}

\subsubsection{Escalation Routing}
\sectionstub{Who is entitled to answer a request, where it appears, and
what happens when nobody answers in time.}
% MATERIAL: §3.5.83.
% DIVERGENCE from Chapter 1 (report/DIVERGENCES.md §1.5): §1.6 puts the
%   timeout with Root and adds a per-user toggle. What was built sets
%   escalation, posture and timeout per AGENT, by an Administrator, and
%   has no per-user toggle.

\subsubsection{Persistent Approvals}
\sectionstub{What the ``always allow'' answer promises, what it does in
practice, and the rule request it files for an administrator.
Figure~\ref{fig:gov-alwaysallow}.}
% FIGURE F23 -> \ref{fig:gov-alwaysallow} (prints as Figure 3.11)
% MATERIAL: §3.5.81, §3.5.94.

\subsection{Runtime Management}
\label{sec:gov-runs}
\sectionstub{How running work is represented, bounded and governed,
including when the agent executes outside this process.}

\subsubsection{Task and Slot Model}
\sectionstub{The row an operator sees and the concurrency slot behind
it, and why the two are released at different moments.
Figure~\ref{fig:gov-taskslot}.}
% FIGURE F24 -> \ref{fig:gov-taskslot} (prints as Figure 3.12)
% MATERIAL: §3.5.82.

\subsubsection{Secondary Runtime Governance}
\sectionstub{The in-process path and the native harness, and the
two-layer permission that governs the second.
Figures~\ref{fig:gov-twopaths} and~\ref{fig:gov-codexgates}.}
% FIGURE F8  -> \ref{fig:gov-twopaths}   (prints as Figure 3.6)
% FIGURE F21 -> \ref{fig:gov-codexgates} (prints as Figure 3.9)
% MATERIAL: §3.5.15, §3.5.62.

\subsection{Tenancy and Agent Registry}
\label{sec:gov-tenant}
\sectionstub{One organization per installation, the storage that belongs
to it, and the registry that gives the layer a noun for an agent.
Figure~\ref{fig:gov-tenant}.}
% FIGURE F19 -> \ref{fig:gov-tenant} (prints as Figure 3.8)
% MATERIAL: §3.5.31, §3.5.33, §3.5.47, §3.5.89.

\subsection{Agent Lifecycle}
\label{sec:gov-lifecycle}
\sectionstub{Registration as the precondition for being allowed
anything; ownership and re-owning; and the two deletions, with what each
one leaves behind on the host.}
% MATERIAL: §3.5.33, §3.5.90, §3.5.91, §3.5.94.

\subsection{Management Interface}
\label{sec:gov-dashboard}
\sectionstub{The HTTP control plane and the dashboard behind it, what
its sections are for, and the rule that a control must never promise
what the gate cannot keep.}
% MATERIAL: §3.5.13, §3.5.37, §3.5.46, §3.5.80.

\subsection{System Security}
\label{sec:gov-security}
\sectionstub{The threats this design answers, and the three it does not:
the check-then-open window, a denied file still reachable on the native
harness, and the fact that the layer governs what an agent does rather
than why it does it. Figure~\ref{fig:gov-toctou}.}
% FIGURE F11 -> \ref{fig:gov-toctou} (prints as Figure 3.7)
% MATERIAL: §3.5.7, §3.5.61; mg/HANDOFF.md §7.
% CAUTION: do not write that the layer prevents prompt injection. It
%   contains the damage when persuasion succeeds. Chapter 2 sets this
%   up, so this section should answer Chapter 2 directly.

% =====================================================================
\section{Summary}
\label{sec:gov-designsummary}

\sectionstub{What the chapter established, the one requirement left
partially met, and the handover to Chapter 4, which tests the design
against Section 1.3 as its benchmark.}
. The \sectionstub definition has
% to go in the preamble; it is repeated at the top of the body so the
% file compiles on its own, and can be deleted once the chapter is done.
%
% The chapter is titled "System Design" because that is what §1.8
% (Document Organization) promises the reader.
%
% ---------------------------------------------------------------------
% HOUSE STYLE FOR THIS CHAPTER
% ---------------------------------------------------------------------
% 1. HEADINGS ARE NOUN PHRASES, two to four words, title case. Not
%    sentences and not questions, following the two model reports in
%    Documentation/GradProj/, whose Chapter 3 headings read "Packet
%    Handling", "Voting Process", "System Architecture", "Smart
%    Contract", "Donor Setup", "System Security". The argument goes in
%    the prose, not the heading.
% 2. AMERICAN SPELLING, because Chapters 1 and 2 already use it:
%    "prioritize", "utilizes", "Data Sanitization", "organizations",
%    "behavior". So: canonicalization, sanitization, organization,
%    authorization, behavior. The repository's own notes are British
%    English; do not carry that into the report.
% 3. The section spine (requirements, analysis of requirements, analysis
%    of constraints, different design approaches, then Developed Design
%    last before the summary) is the spine both model reports use, and
%    the one CHAPTER3-MATERIAL.md is keyed to.
%
% ---------------------------------------------------------------------
% WHERE THE BUILT DESIGN DIFFERS FROM CHAPTERS 1 AND 2
% ---------------------------------------------------------------------
% It is normal and expected that the final implementation differs from
% the preliminary design. What matters is that Chapter 3 states each
% difference openly AND explains why the change was made. The register
% of them, with the justification for each, is
% docs-notes/report/DIVERGENCES.md. Sections that owe an explanation
% carry a % DIVERGENCE comment pointing at it. Add to that file when a
% new one is found, and add the reasoning, not just the fact.
%
% ---------------------------------------------------------------------
% FIGURES: how they are referenced, and where to find them
% ---------------------------------------------------------------------
% The drawings are NOT in this file. They live in docs-notes/FIGURES.md,
% which holds each figure three ways (prose, Mermaid, TikZ) with the
% reasoning behind it. That file is the source of truth.
%
% Prose cites a figure by the label FIGURES.md already gives it, so the
% two files cannot drift:
%
%     ...as shown in Figure~\ref{fig:gov-architecture}.
%
% Every Chapter 3 and 4 figure label is prefixed "fig:gov-". The prefix
% is deliberate: Chapter 1 already uses \label{fig:architecture} for the
% preliminary design figure, and without the prefix our F1 would have
% silently collided with it. Keep the prefix for anything new. Tables
% take "tab:gov-" for the same reason.
%
% TO SEE a figure:  open docs-notes/figures/report-figures.pdf at the
%                   page in the table below.
% TO GET its TikZ:  copy it from docs-notes/FIGURES.md under
%                   "## F1: ..." -> "### TikZ form", or lift the whole
%                   figure environment out of the generated file
%                   docs-notes/figures/report-figures.tex.
% TO REBUILD both:  node docs-notes/figures/build-figures.mjs
%                   then run pdflatex on report-figures.tex.
%
% All fourteen compile clean and have been checked by eye (2026-09-20).
%
%   Prints as | FIGURES.md | Cite with                    | PDF page
%   ----------+------------+------------------------------+---------
%   Fig 3.1   | F1         | \ref{fig:gov-architecture}   |  1
%   Fig 3.2   | F2         | \ref{fig:gov-rbac}           |  2
%   Fig 3.3   | F3         | \ref{fig:gov-decision}       |  3
%   Fig 3.4   | F5         | \ref{fig:gov-pathnorm}       |  4
%   Fig 3.5   | F6         | \ref{fig:gov-promptpath}     |  5
%   Fig 3.6   | F8         | \ref{fig:gov-twopaths}       |  6
%   Fig 3.7   | F11        | \ref{fig:gov-toctou}         |  7
%   Fig 3.8   | F19        | \ref{fig:gov-tenant}         |  8
%   Fig 3.9   | F21        | \ref{fig:gov-codexgates}     |  9
%   Fig 3.10  | F22        | \ref{fig:gov-foldergrant}    | 10
%   Fig 3.11  | F23        | \ref{fig:gov-alwaysallow}    | 11
%   Fig 3.12  | F24        | \ref{fig:gov-taskslot}       | 12
%   Fig 4.1   | F14        | \ref{fig:gov-coverage}       | 13  (Chapter 4)
%   Fig 4.2   | F17        | \ref{fig:gov-defectage}      | 14  (Chapter 4)
%
% main.tex already has tikz, the \usetikzlibrary line, pgfplots with
% compat=1.18, and graphicx (via style-PSUT-BSc-Eng). It still needs the
% shared \tikzset style block from FIGURES.md -> "One shared style
% block", pasted after \pgfplotsset{compat=1.18}. Without it the figures
% compile but lose every box, arrow and label style.
% =====================================================================

% Marks an unwritten section. Put this in the main.tex preamble; delete
% it, and every use of it, once the chapter is written.
\providecommand{\sectionstub}[1]{\emph{[Outline: #1]}\par}

\chapter{System Design}
\label{ch:design}

The preceding chapter established that OS-level Agents introduce a kind of
vulnerability that is cognitive rather than deterministic, and that what an
attacker often manipulates is the reasoning of the agent itself. Any control
that lives inside the model, or that relies on the model's own judgment to
enforce it, therefore inherits the weakness it is meant to correct. The design
in this chapter, briefly described in previous sections, starts with the
assumption that the agent is untrusted, and a deterministic policy check is
placed on the path that every tool call by an agent takes. This required that
the governance layer is built as a hard fork compiled into the OpenClaw core
rather than as a plugin installed alongside it. This commits the project to an
entire upstream codebase, including native client applications that this work
does not modify.

On that foundation the layer adds five capabilities to the host: a default-deny
policy gate on the tool-call path, a tamper-evident audit ledger, an emergency
kill-switch that measures the time to a confirmed stop, four tiers of named
human accounts with individual sessions and levels of authority, and a
management interface through which each of these is operated. The chapter moves
from requirement to finished system. Section~\ref{sec:gov-requirements} briefly
restates the nine design requirements of Section~1.3.
Section~\ref{sec:gov-reqanalysis} then analyzes them, describing the status the
system has reached in relation to each one.
Section~\ref{sec:gov-constraints} does the same for the design constraints of
Section~1.4, and Section~\ref{sec:gov-standards} describes how the engineering
standards of Section~1.5 are applied. Section~\ref{sec:gov-approaches} presents
the alternatives that were seriously considered and the reason each was
rejected, and Section~\ref{sec:gov-developed} describes the system as built, one
design area at a time, ending with what the system does not defend against; in
other words, accepted limitations. Where the developed design departs from the
preliminary design of Chapter~1, the departure is stated where it arises and the
reason is given with it.

% =====================================================================
\section{Design Requirements}
\label{sec:gov-requirements}

As stated in Chapter 1, the design shall achieve the following requirements:

\begin{enumerate}
    \item The system shall be implemented using the Node.js runtime (version 18 or higher) and developed primarily in TypeScript with static type checking enabled.
    \item The system shall provide a secure web-based dashboard that enables administrators to configure customized privilege policies, monitor active autonomous agent sessions, and manage system operations using role-based access control.
    \item The system shall enforce a default-deny, policy-based security model that restricts autonomous agent access to operating system resources, including file system paths, process execution, and network communication, based on administrator-defined capabilities.
    \item The system shall support customized, fine-grained privileges for autonomous agents, including path-level file access, command allowlisting, network allowlisting, and time-limited permissions.
    \item The system shall continuously monitor autonomous agent activities and shall record 100\% of agent actions, policy decisions, and administrative approvals in an auditable logging system.
    \item The system shall implement tamper-evident audit logging mechanisms to ensure the integrity and traceability of all recorded agent and administrative actions.
    \item The system shall provide real-time administrative control capabilities, including the ability to suspend or terminate an autonomous agent session within a maximum response time of one second.
    \item The system shall prevent sensitive data (such as secrets or credentials) from being written in plaintext to log files.
    \item The system shall be deployable on a Linux-based operating system using open-source software components only.
\end{enumerate}


% =====================================================================
\section{Analysis of Design Requirements}
\label{sec:gov-reqanalysis}

Section~\ref{sec:gov-requirements} states the requirements as specified. This
section examines what each one demands of the design when it is read strictly,
and records the status the implementation has reached against it. The
distinction is material, because each requirement is a single sentence of
specification and several of them constrain the design considerably more
tightly than their length suggests. Two turn on the interpretation of a single
phrase, and a third was satisfied only after its initial interpretation was
discarded and the implementation replaced. Table~\ref{tab:gov-reqstatus}
summarizes the outcome and indicates where each requirement is discussed. All
nine requirements are met. The ninth, deployment on a Linux host, was the last
to close, and the evidence that closed it is given in
Section~\ref{sec:gov-remaining-reqs}.

\begin{table}[h]
\centering
\caption{Design requirement status}
\label{tab:gov-reqstatus}
\renewcommand{\arraystretch}{1.3}
\small
\begin{tabular}{@{}c p{5.7cm} l l@{}}
\toprule
\textbf{\#} & \textbf{Requirement} & \textbf{Status} & \textbf{Discussed in} \\
\midrule
1 & Node.js and TypeScript, with static type checking & Met & \S\ref{sec:gov-remaining-reqs} \\
2 & Secure web dashboard with role-based access control & Met & \S\ref{sec:gov-dashboard} \\
3 & Default-deny model over paths, execution and network & Met & \S\ref{sec:gov-enforcement} \\
4 & Fine-grained, time-limited privileges & Met & \S\ref{sec:gov-gate} \\
5 & Record 100\% of actions, decisions and approvals & Met & \S\ref{sec:gov-coverage} \\
6 & Tamper-evident audit logging & Met & \S\ref{sec:gov-ledger} \\
7 & Suspend or terminate a session within one second & Met & \S\ref{sec:gov-latency} \\
8 & No plaintext secrets in log files & Met & \S\ref{sec:gov-credentials} \\
9 & Deployable on Linux, open-source components only & Met & \S\ref{sec:gov-remaining-reqs} \\
\bottomrule
\end{tabular}
\end{table}

Four requirements are treated individually below, because in each case a strict
reading either altered the design or altered what could honestly be claimed of
it. Requirement 3 turned on the position of the check rather than on the
decision it reaches. Requirement 5 turned on the scope of \emph{every action},
and its first implementation was replaced. Requirement 7 turned on the quantity
being measured, the dispatch of a stop signal or the stop itself. Requirement 8
was satisfied by reusing an existing component rather than implementing a second
one, which is a design decision and is presented as such. The remaining five
requirements are addressed together in Section~\ref{sec:gov-remaining-reqs}.

\subsection{Enforcement Point}
\label{sec:gov-enforcement}
\sectionstub{Requirement 3. Why a default-deny gate has to sit on the
host's own tool-call path rather than beside it, and what that costs in
coverage.}

\subsection{Audit Coverage}
\label{sec:gov-coverage}
\sectionstub{Requirement 5. What recording one hundred percent of agent
actions was first read as meaning, why that reading was abandoned, and
the distinct decision that lets an unevaluated call be recorded as such
rather than folded in with the allowed ones.}

\subsection{Termination Latency}
\label{sec:gov-latency}
\sectionstub{Requirement 7. The difference between the stop signal being
sent and the stop being confirmed, and which of the two the requirement
is read as claiming.}

\subsection{Credential Protection}
\label{sec:gov-credentials}
\sectionstub{Requirement 8. Met by redacting at the ledger boundary with
the host's own redaction rather than by reimplementing it.}

\subsection{Remaining Requirements}
\label{sec:gov-remaining-reqs}
\sectionstub{Requirements 1, 2, 4, 6 and 9 in brief: what each asked,
and what satisfying it implied for the design.}

% =====================================================================
\section{Analysis of Design Constraints}
\label{sec:gov-constraints}

\sectionstub{The three categories of Section 1.4, and what each one
ruled out of the design.}

% MATERIAL: CHAPTER3-MATERIAL.md "→ 3.3 Analysis of Design Constraints".

\subsection{Economic Constraints}
\sectionstub{The 350 JOD ceiling and the open-source-only stack, met in
part by adding no new dependencies at all, which is a claim any reader
can check.}

\subsection{Manufacturability and Sustainability}
\sectionstub{What the layer adds to resident memory and to startup time
on an 8\,GB server, and the startup budget that was set and defended
because of this constraint.}

\subsection{Ethical and Professional Responsibility}
\sectionstub{The design can subtract capability from an agent and never
add it, and why that asymmetry was made structural rather than left to
care.}

\subsection{Accepted Limitations}
\sectionstub{Constraints the design accepted rather than solved:
English-only operator text, and development on Windows against a Linux
deployment target.}

% =====================================================================
\section{Discussion of Engineering Standards}
\label{sec:gov-standards}

\sectionstub{How the design applies the two standards named in Section
1.5, rather than restating what those standards are.}

% NOTE: the one section with no material behind it in the repository.
%   Write it from §1.5's wording plus the design as built.

\subsection{ISO/IEC 27001:2022}
\sectionstub{Policy-driven access control mapped onto the default-deny
gate and its rule language; systematic governance of rights onto the
four tiers and ownership; audit logging onto the hash-chained ledger;
and risk treatment onto the baseline policy and the core denials that
cannot be switched off at runtime.}

\subsection{OWASP Secure Coding Practices}
\sectionstub{Session management and authentication mapped onto the two
gates on every governance request; error handling onto refusals that
explain themselves without leaking; data in transit onto keeping the
dashboard off the public internet behind an SSH tunnel; and input
handling onto rule validation and path canonicalization.}

% =====================================================================
\section{Different Design Approaches}
\label{sec:gov-approaches}

\sectionstub{The alternatives that were genuinely weighed, each stated
as the options and then the choice, with the reason the others lost.}

% MATERIAL: CHAPTER3-MATERIAL.md "→ 3.4" holds two of these in full; the
%   rest are recoverable from the §3.5 sections named under each heading.

\subsection{Integration Strategy}
\sectionstub{A fork compiled into the host's core, or a plugin beside
it. Why a plugin cannot make a gate unavoidable.}
% MATERIAL: §3.5.2b.

\subsection{Gate Placement}
\sectionstub{Where in the tool-call path the check belongs, and what
each candidate position would and would not have seen.}
% MATERIAL: §3.5.1, §3.5.25.

\subsection{Host Interception}
\sectionstub{How the host is told that the gate must be consulted: four
options, one chosen.}
% MATERIAL: "→ 3.4" §3.4.y, §3.5.15.

\subsection{Path Representation}
\sectionstub{Which form a file path takes when a rule is matched against
it: three options, one chosen.}
% MATERIAL: "→ 3.4" §3.4.x, §3.5.8.

\subsection{Log Integrity}
\sectionstub{A hash chain weighed against the alternatives, and what
each would have cost in complexity and in trust.}
% MATERIAL: §3.5.3, §3.5.9.

\subsection{Operator Identity}
\sectionstub{One shared Gateway credential, or a second and independent
set of named accounts with their own sessions.}
% MATERIAL: §3.5.6, §3.5.19.

\subsection{Management Surfaces}
\sectionstub{Three surfaces were built and one was removed, and why
subtraction counts as a design decision rather than a retreat.}
% MATERIAL: §3.5.77.

\subsection{Tenancy Model}
\sectionstub{How many organizations one installation holds. One, decided
deliberately, with the structure for several deliberately kept.}
% MATERIAL: §3.5.89.

% =====================================================================
\section{Developed Design}
\label{sec:gov-developed}

\sectionstub{The system as built, taken one design area at a time. Each
part is introduced by what it is for, then how it works, then the one
decision inside it that a reader would otherwise stop and question.}

\subsection{System Architecture}
\label{sec:gov-arch}
\sectionstub{The layer inside the Gateway process; the two doors into
it, one for the operator and one for agent activity; and the files it
owns on disk. Figure~\ref{fig:gov-architecture}.}
% FIGURE F1 -> \ref{fig:gov-architecture} (prints as Figure 3.1)
% MATERIAL: §3.5.1, §3.5.2, §3.5.2b.

\subsection{Policy Engine}
\label{sec:gov-gate}
\sectionstub{The component that turns a rule set and an attempted action
into a decision.}

\subsubsection{Rule Model}
\sectionstub{Tiers, effects, access directions, resource kinds and
expiry: the vocabulary a rule is written in.}
% MATERIAL: §3.5.4, §3.5.12, §3.5.20; docs-notes/PERMISSION-SPEC.md.

\subsubsection{Evaluation Order}
\sectionstub{Why denials are read before allowances, and what that
guarantees. Figure~\ref{fig:gov-decision}.}
% FIGURE F3 -> \ref{fig:gov-decision} (prints as Figure 3.3)
% MATERIAL: §3.5.5, §3.5.25.
% DIVERGENCE from Chapter 1 (report/DIVERGENCES.md §1.3): §1.6 says
%   conflicts are resolved by preferring the earlier rule. That is true
%   of the conflict DETECTOR but not of evaluation, where a denial beats
%   an allowance whatever the order. Complete the picture here.

\subsubsection{Path Canonicalization}
\sectionstub{Turning what an agent wrote into the path it resolves to on
disk, before any rule is matched against it.
Figure~\ref{fig:gov-pathnorm}.}
% FIGURE F5 -> \ref{fig:gov-pathnorm} (prints as Figure 3.4)
% MATERIAL: §3.5.8, §3.5.29.

\subsubsection{Baseline Policy}
\sectionstub{What an installation ships with: core denials that cannot
be edited at runtime, and the allowances that make an agent useful
before anyone has written a rule.}
% MATERIAL: docs-notes/BASELINE-RULES.md.

\subsubsection{Folder Grants}
\sectionstub{Granting a directory with exceptions carved out of it, and
why the denial is written first. Figure~\ref{fig:gov-foldergrant}.}
% FIGURE F22 -> \ref{fig:gov-foldergrant} (prints as Figure 3.10)
% MATERIAL: §3.5.66.

\subsection{Audit Ledger}
\label{sec:gov-ledger}
\sectionstub{The append-only record of what was attempted, what was
decided and who decided it.}

\subsubsection{Entry Structure}
\sectionstub{What a single entry holds, including the raw model intent
that produced the action.}
% MATERIAL: §3.5.3, §3.5.59.

\subsubsection{Hash Chaining and Verification}
\sectionstub{How each entry is bound to the one before it, and the
independent verifier that detects any later edit.}
% MATERIAL: §3.5.3; scripts/verify-ledger.mjs.
% DIVERGENCE from Chapter 2 (report/DIVERGENCES.md §1.1): Chapter 2
%   describes a Merkle tree; a keyed HMAC-SHA256 chain was built. State
%   the difference AND why: the tree's advantage is proving one entry
%   without reading the log, which no party here needs, and the keying
%   is a strengthening over Chapter 2's plain SHA-256.

\subsubsection{Administrative Logging}
\sectionstub{Why administrative actions are recorded in the same chain
as agent actions, and what each such entry must name.}
% MATERIAL: §3.5.9, §3.5.63.

\subsubsection{Data Sanitization}
\sectionstub{What an audit trail is allowed not to see, and where
redaction is applied so that it cannot be bypassed.}
% MATERIAL: §3.5.28, §3.5.60.
% DIVERGENCE from Chapter 2 (report/DIVERGENCES.md §1.2): Chapter 2
%   promises regex and entropy-based masking; what was built reuses the
%   host's own redactor at the ledger boundary. State the difference AND
%   why: one tested redactor rather than two to maintain, and placement
%   at the single write boundary beats cleverness applied per caller.

\subsection{Access Control Model}
\label{sec:gov-rbac}
\sectionstub{Who the operators are, what each of them may do, and how
the system decides. Figure~\ref{fig:gov-rbac}.}
% FIGURE F2 -> \ref{fig:gov-rbac} (prints as Figure 3.2)

\subsubsection{Role Hierarchy}
\sectionstub{The four tiers, what each inherits from the one below, and
the capability that deliberately does not sit at the top.}
% MATERIAL: §3.5.4, §3.5.24; docs-notes/ROLE-MODEL.md.

\subsubsection{Two-Gate Authentication}
\sectionstub{The Gateway credential, and then the governance account
session and its role, on every request.}
% MATERIAL: §3.5.6, §3.5.19.

\subsubsection{Ownership and Assignment}
\sectionstub{What managing an agent resolves to in code: the two scope
questions each tier answers independently.}
% MATERIAL: docs-notes/ROLE-MODEL.md, §3.5.22, §3.5.30.

\subsection{Prompt Execution Path}
\label{sec:gov-promptpath}
\sectionstub{A User prompting an agent they hold: the checks made, and
the record written before the run rather than after it.
Figure~\ref{fig:gov-promptpath}.}
% FIGURE F6 -> \ref{fig:gov-promptpath} (prints as Figure 3.5)
% MATERIAL: §3.5.11, §3.5.17.

\subsection{Session Control}
\label{sec:gov-killswitch}
\sectionstub{The two ways an operator intervenes in a running agent, and
what each one is able to reach.}

\subsubsection{Agent Lockdown}
\sectionstub{What a lockdown covers, including work the locked agent had
already started, and why it outranks any allowing rule.}
% MATERIAL: §3.5.38, §3.5.16.

\subsubsection{Kill Switch}
\sectionstub{Stopping a run, and why the design reports a confirmed stop
rather than a signal that was sent.}
% MATERIAL: §3.5.10, §3.5.21.
% DIVERGENCE from Chapter 1 (report/DIVERGENCES.md §1.4): §1.6 promises
%   a sub-second kill switch. The signal is milliseconds; the CONFIRMED
%   stop measured 2.2 s and 2.8 s on the laptop. Give both numbers and
%   say the design measures something stricter than Chapter 1 did.

\subsection{Human-in-the-Loop Approvals}
\label{sec:gov-approvals}
\sectionstub{What happens when the policy neither allows nor denies an
action and a person has to decide.}

\subsubsection{Escalation Routing}
\sectionstub{Who is entitled to answer a request, where it appears, and
what happens when nobody answers in time.}
% MATERIAL: §3.5.83.
% DIVERGENCE from Chapter 1 (report/DIVERGENCES.md §1.5): §1.6 puts the
%   timeout with Root and adds a per-user toggle. What was built sets
%   escalation, posture and timeout per AGENT, by an Administrator, and
%   has no per-user toggle.

\subsubsection{Persistent Approvals}
\sectionstub{What the ``always allow'' answer promises, what it does in
practice, and the rule request it files for an administrator.
Figure~\ref{fig:gov-alwaysallow}.}
% FIGURE F23 -> \ref{fig:gov-alwaysallow} (prints as Figure 3.11)
% MATERIAL: §3.5.81, §3.5.94.

\subsection{Runtime Management}
\label{sec:gov-runs}
\sectionstub{How running work is represented, bounded and governed,
including when the agent executes outside this process.}

\subsubsection{Task and Slot Model}
\sectionstub{The row an operator sees and the concurrency slot behind
it, and why the two are released at different moments.
Figure~\ref{fig:gov-taskslot}.}
% FIGURE F24 -> \ref{fig:gov-taskslot} (prints as Figure 3.12)
% MATERIAL: §3.5.82.

\subsubsection{Secondary Runtime Governance}
\sectionstub{The in-process path and the native harness, and the
two-layer permission that governs the second.
Figures~\ref{fig:gov-twopaths} and~\ref{fig:gov-codexgates}.}
% FIGURE F8  -> \ref{fig:gov-twopaths}   (prints as Figure 3.6)
% FIGURE F21 -> \ref{fig:gov-codexgates} (prints as Figure 3.9)
% MATERIAL: §3.5.15, §3.5.62.

\subsection{Tenancy and Agent Registry}
\label{sec:gov-tenant}
\sectionstub{One organization per installation, the storage that belongs
to it, and the registry that gives the layer a noun for an agent.
Figure~\ref{fig:gov-tenant}.}
% FIGURE F19 -> \ref{fig:gov-tenant} (prints as Figure 3.8)
% MATERIAL: §3.5.31, §3.5.33, §3.5.47, §3.5.89.

\subsection{Agent Lifecycle}
\label{sec:gov-lifecycle}
\sectionstub{Registration as the precondition for being allowed
anything; ownership and re-owning; and the two deletions, with what each
one leaves behind on the host.}
% MATERIAL: §3.5.33, §3.5.90, §3.5.91, §3.5.94.

\subsection{Management Interface}
\label{sec:gov-dashboard}
\sectionstub{The HTTP control plane and the dashboard behind it, what
its sections are for, and the rule that a control must never promise
what the gate cannot keep.}
% MATERIAL: §3.5.13, §3.5.37, §3.5.46, §3.5.80.

\subsection{System Security}
\label{sec:gov-security}
\sectionstub{The threats this design answers, and the three it does not:
the check-then-open window, a denied file still reachable on the native
harness, and the fact that the layer governs what an agent does rather
than why it does it. Figure~\ref{fig:gov-toctou}.}
% FIGURE F11 -> \ref{fig:gov-toctou} (prints as Figure 3.7)
% MATERIAL: §3.5.7, §3.5.61; mg/HANDOFF.md §7.
% CAUTION: do not write that the layer prevents prompt injection. It
%   contains the damage when persuasion succeeds. Chapter 2 sets this
%   up, so this section should answer Chapter 2 directly.

% =====================================================================
\section{Summary}
\label{sec:gov-designsummary}

\sectionstub{What the chapter established, the one requirement left
partially met, and the handover to Chapter 4, which tests the design
against Section 1.3 as its benchmark.}
. The \sectionstub definition has
% to go in the preamble; it is repeated at the top of the body so the
% file compiles on its own, and can be deleted once the chapter is done.
%
% The chapter is titled "System Design" because that is what §1.8
% (Document Organization) promises the reader.
%
% ---------------------------------------------------------------------
% HOUSE STYLE FOR THIS CHAPTER
% ---------------------------------------------------------------------
% 1. HEADINGS ARE NOUN PHRASES, two to four words, title case. Not
%    sentences and not questions, following the two model reports in
%    Documentation/GradProj/, whose Chapter 3 headings read "Packet
%    Handling", "Voting Process", "System Architecture", "Smart
%    Contract", "Donor Setup", "System Security". The argument goes in
%    the prose, not the heading.
% 2. AMERICAN SPELLING, because Chapters 1 and 2 already use it:
%    "prioritize", "utilizes", "Data Sanitization", "organizations",
%    "behavior". So: canonicalization, sanitization, organization,
%    authorization, behavior. The repository's own notes are British
%    English; do not carry that into the report.
% 3. The section spine (requirements, analysis of requirements, analysis
%    of constraints, different design approaches, then Developed Design
%    last before the summary) is the spine both model reports use, and
%    the one CHAPTER3-MATERIAL.md is keyed to.
%
% ---------------------------------------------------------------------
% WHERE THE BUILT DESIGN DIFFERS FROM CHAPTERS 1 AND 2
% ---------------------------------------------------------------------
% It is normal and expected that the final implementation differs from
% the preliminary design. What matters is that Chapter 3 states each
% difference openly AND explains why the change was made. The register
% of them, with the justification for each, is
% docs-notes/report/DIVERGENCES.md. Sections that owe an explanation
% carry a % DIVERGENCE comment pointing at it. Add to that file when a
% new one is found, and add the reasoning, not just the fact.
%
% ---------------------------------------------------------------------
% FIGURES: how they are referenced, and where to find them
% ---------------------------------------------------------------------
% The drawings are NOT in this file. They live in docs-notes/FIGURES.md,
% which holds each figure three ways (prose, Mermaid, TikZ) with the
% reasoning behind it. That file is the source of truth.
%
% Prose cites a figure by the label FIGURES.md already gives it, so the
% two files cannot drift:
%
%     ...as shown in Figure~\ref{fig:gov-architecture}.
%
% Every Chapter 3 and 4 figure label is prefixed "fig:gov-". The prefix
% is deliberate: Chapter 1 already uses \label{fig:architecture} for the
% preliminary design figure, and without the prefix our F1 would have
% silently collided with it. Keep the prefix for anything new. Tables
% take "tab:gov-" for the same reason.
%
% TO SEE a figure:  open docs-notes/figures/report-figures.pdf at the
%                   page in the table below.
% TO GET its TikZ:  copy it from docs-notes/FIGURES.md under
%                   "## F1: ..." -> "### TikZ form", or lift the whole
%                   figure environment out of the generated file
%                   docs-notes/figures/report-figures.tex.
% TO REBUILD both:  node docs-notes/figures/build-figures.mjs
%                   then run pdflatex on report-figures.tex.
%
% All fourteen compile clean and have been checked by eye (2026-09-20).
%
%   Prints as | FIGURES.md | Cite with                    | PDF page
%   ----------+------------+------------------------------+---------
%   Fig 3.1   | F1         | \ref{fig:gov-architecture}   |  1
%   Fig 3.2   | F2         | \ref{fig:gov-rbac}           |  2
%   Fig 3.3   | F3         | \ref{fig:gov-decision}       |  3
%   Fig 3.4   | F5         | \ref{fig:gov-pathnorm}       |  4
%   Fig 3.5   | F6         | \ref{fig:gov-promptpath}     |  5
%   Fig 3.6   | F8         | \ref{fig:gov-twopaths}       |  6
%   Fig 3.7   | F11        | \ref{fig:gov-toctou}         |  7
%   Fig 3.8   | F19        | \ref{fig:gov-tenant}         |  8
%   Fig 3.9   | F21        | \ref{fig:gov-codexgates}     |  9
%   Fig 3.10  | F22        | \ref{fig:gov-foldergrant}    | 10
%   Fig 3.11  | F23        | \ref{fig:gov-alwaysallow}    | 11
%   Fig 3.12  | F24        | \ref{fig:gov-taskslot}       | 12
%   Fig 4.1   | F14        | \ref{fig:gov-coverage}       | 13  (Chapter 4)
%   Fig 4.2   | F17        | \ref{fig:gov-defectage}      | 14  (Chapter 4)
%
% main.tex already has tikz, the \usetikzlibrary line, pgfplots with
% compat=1.18, and graphicx (via style-PSUT-BSc-Eng). It still needs the
% shared \tikzset style block from FIGURES.md -> "One shared style
% block", pasted after \pgfplotsset{compat=1.18}. Without it the figures
% compile but lose every box, arrow and label style.
% =====================================================================

% Marks an unwritten section. Put this in the main.tex preamble; delete
% it, and every use of it, once the chapter is written.
\providecommand{\sectionstub}[1]{\emph{[Outline: #1]}\par}

\chapter{System Design}
\label{ch:design}

The preceding chapter established that OS-level Agents introduce a kind of
vulnerability that is cognitive rather than deterministic, and that what an
attacker often manipulates is the reasoning of the agent itself. Any control
that lives inside the model, or that relies on the model's own judgment to
enforce it, therefore inherits the weakness it is meant to correct. The design
in this chapter, briefly described in previous sections, starts with the
assumption that the agent is untrusted, and a deterministic policy check is
placed on the path that every tool call by an agent takes. This required that
the governance layer is built as a hard fork compiled into the OpenClaw core
rather than as a plugin installed alongside it. This commits the project to an
entire upstream codebase, including native client applications that this work
does not modify.

On that foundation the layer adds five capabilities to the host: a default-deny
policy gate on the tool-call path, a tamper-evident audit ledger, an emergency
kill-switch that measures the time to a confirmed stop, four tiers of named
human accounts with individual sessions and levels of authority, and a
management interface through which each of these is operated. The chapter moves
from requirement to finished system. Section~\ref{sec:gov-requirements} briefly
restates the nine design requirements of Section~1.3.
Section~\ref{sec:gov-reqanalysis} then analyzes them, describing the status the
system has reached in relation to each one.
Section~\ref{sec:gov-constraints} does the same for the design constraints of
Section~1.4, and Section~\ref{sec:gov-standards} describes how the engineering
standards of Section~1.5 are applied. Section~\ref{sec:gov-approaches} presents
the alternatives that were seriously considered and the reason each was
rejected, and Section~\ref{sec:gov-developed} describes the system as built, one
design area at a time, ending with what the system does not defend against; in
other words, accepted limitations. Where the developed design departs from the
preliminary design of Chapter~1, the departure is stated where it arises and the
reason is given with it.

% =====================================================================
\section{Design Requirements}
\label{sec:gov-requirements}

As stated in Chapter 1, the design shall achieve the following requirements:

\begin{enumerate}
    \item The system shall be implemented using the Node.js runtime (version 18 or higher) and developed primarily in TypeScript with static type checking enabled.
    \item The system shall provide a secure web-based dashboard that enables administrators to configure customized privilege policies, monitor active autonomous agent sessions, and manage system operations using role-based access control.
    \item The system shall enforce a default-deny, policy-based security model that restricts autonomous agent access to operating system resources, including file system paths, process execution, and network communication, based on administrator-defined capabilities.
    \item The system shall support customized, fine-grained privileges for autonomous agents, including path-level file access, command allowlisting, network allowlisting, and time-limited permissions.
    \item The system shall continuously monitor autonomous agent activities and shall record 100\% of agent actions, policy decisions, and administrative approvals in an auditable logging system.
    \item The system shall implement tamper-evident audit logging mechanisms to ensure the integrity and traceability of all recorded agent and administrative actions.
    \item The system shall provide real-time administrative control capabilities, including the ability to suspend or terminate an autonomous agent session within a maximum response time of one second.
    \item The system shall prevent sensitive data (such as secrets or credentials) from being written in plaintext to log files.
    \item The system shall be deployable on a Linux-based operating system using open-source software components only.
\end{enumerate}


% =====================================================================
\section{Analysis of Design Requirements}
\label{sec:gov-reqanalysis}

Section~\ref{sec:gov-requirements} states the requirements as specified. This
section examines what each one demands of the design when it is read strictly,
and records the status the implementation has reached against it. The
distinction is material, because each requirement is a single sentence of
specification and several of them constrain the design considerably more
tightly than their length suggests. Two turn on the interpretation of a single
phrase, and a third was satisfied only after its initial interpretation was
discarded and the implementation replaced. Table~\ref{tab:gov-reqstatus}
summarizes the outcome and indicates where each requirement is discussed. All
nine requirements are met. The ninth, deployment on a Linux host, was the last
to close, and the evidence that closed it is given in
Section~\ref{sec:gov-remaining-reqs}.

\begin{table}[h]
\centering
\caption{Design requirement status}
\label{tab:gov-reqstatus}
\renewcommand{\arraystretch}{1.3}
\small
\begin{tabular}{@{}c p{5.7cm} l l@{}}
\toprule
\textbf{\#} & \textbf{Requirement} & \textbf{Status} & \textbf{Discussed in} \\
\midrule
1 & Node.js and TypeScript, with static type checking & Met & \S\ref{sec:gov-remaining-reqs} \\
2 & Secure web dashboard with role-based access control & Met & \S\ref{sec:gov-dashboard} \\
3 & Default-deny model over paths, execution and network & Met & \S\ref{sec:gov-enforcement} \\
4 & Fine-grained, time-limited privileges & Met & \S\ref{sec:gov-gate} \\
5 & Record 100\% of actions, decisions and approvals & Met & \S\ref{sec:gov-coverage} \\
6 & Tamper-evident audit logging & Met & \S\ref{sec:gov-ledger} \\
7 & Suspend or terminate a session within one second & Met & \S\ref{sec:gov-latency} \\
8 & No plaintext secrets in log files & Met & \S\ref{sec:gov-credentials} \\
9 & Deployable on Linux, open-source components only & Met & \S\ref{sec:gov-remaining-reqs} \\
\bottomrule
\end{tabular}
\end{table}

Four requirements are treated individually below, because in each case a strict
reading either altered the design or altered what could honestly be claimed of
it. Requirement 3 turned on the position of the check rather than on the
decision it reaches. Requirement 5 turned on the scope of \emph{every action},
and its first implementation was replaced. Requirement 7 turned on the quantity
being measured, the dispatch of a stop signal or the stop itself. Requirement 8
was satisfied by reusing an existing component rather than implementing a second
one, which is a design decision and is presented as such. The remaining five
requirements are addressed together in Section~\ref{sec:gov-remaining-reqs}.

\subsection{Enforcement Point}
\label{sec:gov-enforcement}
\sectionstub{Requirement 3. Why a default-deny gate has to sit on the
host's own tool-call path rather than beside it, and what that costs in
coverage.}

\subsection{Audit Coverage}
\label{sec:gov-coverage}
\sectionstub{Requirement 5. What recording one hundred percent of agent
actions was first read as meaning, why that reading was abandoned, and
the distinct decision that lets an unevaluated call be recorded as such
rather than folded in with the allowed ones.}

\subsection{Termination Latency}
\label{sec:gov-latency}
\sectionstub{Requirement 7. The difference between the stop signal being
sent and the stop being confirmed, and which of the two the requirement
is read as claiming.}

\subsection{Credential Protection}
\label{sec:gov-credentials}
\sectionstub{Requirement 8. Met by redacting at the ledger boundary with
the host's own redaction rather than by reimplementing it.}

\subsection{Remaining Requirements}
\label{sec:gov-remaining-reqs}
\sectionstub{Requirements 1, 2, 4, 6 and 9 in brief: what each asked,
and what satisfying it implied for the design.}

% =====================================================================
\section{Analysis of Design Constraints}
\label{sec:gov-constraints}

\sectionstub{The three categories of Section 1.4, and what each one
ruled out of the design.}

% MATERIAL: CHAPTER3-MATERIAL.md "→ 3.3 Analysis of Design Constraints".

\subsection{Economic Constraints}
\sectionstub{The 350 JOD ceiling and the open-source-only stack, met in
part by adding no new dependencies at all, which is a claim any reader
can check.}

\subsection{Manufacturability and Sustainability}
\sectionstub{What the layer adds to resident memory and to startup time
on an 8\,GB server, and the startup budget that was set and defended
because of this constraint.}

\subsection{Ethical and Professional Responsibility}
\sectionstub{The design can subtract capability from an agent and never
add it, and why that asymmetry was made structural rather than left to
care.}

\subsection{Accepted Limitations}
\sectionstub{Constraints the design accepted rather than solved:
English-only operator text, and development on Windows against a Linux
deployment target.}

% =====================================================================
\section{Discussion of Engineering Standards}
\label{sec:gov-standards}

\sectionstub{How the design applies the two standards named in Section
1.5, rather than restating what those standards are.}

% NOTE: the one section with no material behind it in the repository.
%   Write it from §1.5's wording plus the design as built.

\subsection{ISO/IEC 27001:2022}
\sectionstub{Policy-driven access control mapped onto the default-deny
gate and its rule language; systematic governance of rights onto the
four tiers and ownership; audit logging onto the hash-chained ledger;
and risk treatment onto the baseline policy and the core denials that
cannot be switched off at runtime.}

\subsection{OWASP Secure Coding Practices}
\sectionstub{Session management and authentication mapped onto the two
gates on every governance request; error handling onto refusals that
explain themselves without leaking; data in transit onto keeping the
dashboard off the public internet behind an SSH tunnel; and input
handling onto rule validation and path canonicalization.}

% =====================================================================
\section{Different Design Approaches}
\label{sec:gov-approaches}

\sectionstub{The alternatives that were genuinely weighed, each stated
as the options and then the choice, with the reason the others lost.}

% MATERIAL: CHAPTER3-MATERIAL.md "→ 3.4" holds two of these in full; the
%   rest are recoverable from the §3.5 sections named under each heading.

\subsection{Integration Strategy}
\sectionstub{A fork compiled into the host's core, or a plugin beside
it. Why a plugin cannot make a gate unavoidable.}
% MATERIAL: §3.5.2b.

\subsection{Gate Placement}
\sectionstub{Where in the tool-call path the check belongs, and what
each candidate position would and would not have seen.}
% MATERIAL: §3.5.1, §3.5.25.

\subsection{Host Interception}
\sectionstub{How the host is told that the gate must be consulted: four
options, one chosen.}
% MATERIAL: "→ 3.4" §3.4.y, §3.5.15.

\subsection{Path Representation}
\sectionstub{Which form a file path takes when a rule is matched against
it: three options, one chosen.}
% MATERIAL: "→ 3.4" §3.4.x, §3.5.8.

\subsection{Log Integrity}
\sectionstub{A hash chain weighed against the alternatives, and what
each would have cost in complexity and in trust.}
% MATERIAL: §3.5.3, §3.5.9.

\subsection{Operator Identity}
\sectionstub{One shared Gateway credential, or a second and independent
set of named accounts with their own sessions.}
% MATERIAL: §3.5.6, §3.5.19.

\subsection{Management Surfaces}
\sectionstub{Three surfaces were built and one was removed, and why
subtraction counts as a design decision rather than a retreat.}
% MATERIAL: §3.5.77.

\subsection{Tenancy Model}
\sectionstub{How many organizations one installation holds. One, decided
deliberately, with the structure for several deliberately kept.}
% MATERIAL: §3.5.89.

% =====================================================================
\section{Developed Design}
\label{sec:gov-developed}

\sectionstub{The system as built, taken one design area at a time. Each
part is introduced by what it is for, then how it works, then the one
decision inside it that a reader would otherwise stop and question.}

\subsection{System Architecture}
\label{sec:gov-arch}
\sectionstub{The layer inside the Gateway process; the two doors into
it, one for the operator and one for agent activity; and the files it
owns on disk. Figure~\ref{fig:gov-architecture}.}
% FIGURE F1 -> \ref{fig:gov-architecture} (prints as Figure 3.1)
% MATERIAL: §3.5.1, §3.5.2, §3.5.2b.

\subsection{Policy Engine}
\label{sec:gov-gate}
\sectionstub{The component that turns a rule set and an attempted action
into a decision.}

\subsubsection{Rule Model}
\sectionstub{Tiers, effects, access directions, resource kinds and
expiry: the vocabulary a rule is written in.}
% MATERIAL: §3.5.4, §3.5.12, §3.5.20; docs-notes/PERMISSION-SPEC.md.

\subsubsection{Evaluation Order}
\sectionstub{Why denials are read before allowances, and what that
guarantees. Figure~\ref{fig:gov-decision}.}
% FIGURE F3 -> \ref{fig:gov-decision} (prints as Figure 3.3)
% MATERIAL: §3.5.5, §3.5.25.
% DIVERGENCE from Chapter 1 (report/DIVERGENCES.md §1.3): §1.6 says
%   conflicts are resolved by preferring the earlier rule. That is true
%   of the conflict DETECTOR but not of evaluation, where a denial beats
%   an allowance whatever the order. Complete the picture here.

\subsubsection{Path Canonicalization}
\sectionstub{Turning what an agent wrote into the path it resolves to on
disk, before any rule is matched against it.
Figure~\ref{fig:gov-pathnorm}.}
% FIGURE F5 -> \ref{fig:gov-pathnorm} (prints as Figure 3.4)
% MATERIAL: §3.5.8, §3.5.29.

\subsubsection{Baseline Policy}
\sectionstub{What an installation ships with: core denials that cannot
be edited at runtime, and the allowances that make an agent useful
before anyone has written a rule.}
% MATERIAL: docs-notes/BASELINE-RULES.md.

\subsubsection{Folder Grants}
\sectionstub{Granting a directory with exceptions carved out of it, and
why the denial is written first. Figure~\ref{fig:gov-foldergrant}.}
% FIGURE F22 -> \ref{fig:gov-foldergrant} (prints as Figure 3.10)
% MATERIAL: §3.5.66.

\subsection{Audit Ledger}
\label{sec:gov-ledger}
\sectionstub{The append-only record of what was attempted, what was
decided and who decided it.}

\subsubsection{Entry Structure}
\sectionstub{What a single entry holds, including the raw model intent
that produced the action.}
% MATERIAL: §3.5.3, §3.5.59.

\subsubsection{Hash Chaining and Verification}
\sectionstub{How each entry is bound to the one before it, and the
independent verifier that detects any later edit.}
% MATERIAL: §3.5.3; scripts/verify-ledger.mjs.
% DIVERGENCE from Chapter 2 (report/DIVERGENCES.md §1.1): Chapter 2
%   describes a Merkle tree; a keyed HMAC-SHA256 chain was built. State
%   the difference AND why: the tree's advantage is proving one entry
%   without reading the log, which no party here needs, and the keying
%   is a strengthening over Chapter 2's plain SHA-256.

\subsubsection{Administrative Logging}
\sectionstub{Why administrative actions are recorded in the same chain
as agent actions, and what each such entry must name.}
% MATERIAL: §3.5.9, §3.5.63.

\subsubsection{Data Sanitization}
\sectionstub{What an audit trail is allowed not to see, and where
redaction is applied so that it cannot be bypassed.}
% MATERIAL: §3.5.28, §3.5.60.
% DIVERGENCE from Chapter 2 (report/DIVERGENCES.md §1.2): Chapter 2
%   promises regex and entropy-based masking; what was built reuses the
%   host's own redactor at the ledger boundary. State the difference AND
%   why: one tested redactor rather than two to maintain, and placement
%   at the single write boundary beats cleverness applied per caller.

\subsection{Access Control Model}
\label{sec:gov-rbac}
\sectionstub{Who the operators are, what each of them may do, and how
the system decides. Figure~\ref{fig:gov-rbac}.}
% FIGURE F2 -> \ref{fig:gov-rbac} (prints as Figure 3.2)

\subsubsection{Role Hierarchy}
\sectionstub{The four tiers, what each inherits from the one below, and
the capability that deliberately does not sit at the top.}
% MATERIAL: §3.5.4, §3.5.24; docs-notes/ROLE-MODEL.md.

\subsubsection{Two-Gate Authentication}
\sectionstub{The Gateway credential, and then the governance account
session and its role, on every request.}
% MATERIAL: §3.5.6, §3.5.19.

\subsubsection{Ownership and Assignment}
\sectionstub{What managing an agent resolves to in code: the two scope
questions each tier answers independently.}
% MATERIAL: docs-notes/ROLE-MODEL.md, §3.5.22, §3.5.30.

\subsection{Prompt Execution Path}
\label{sec:gov-promptpath}
\sectionstub{A User prompting an agent they hold: the checks made, and
the record written before the run rather than after it.
Figure~\ref{fig:gov-promptpath}.}
% FIGURE F6 -> \ref{fig:gov-promptpath} (prints as Figure 3.5)
% MATERIAL: §3.5.11, §3.5.17.

\subsection{Session Control}
\label{sec:gov-killswitch}
\sectionstub{The two ways an operator intervenes in a running agent, and
what each one is able to reach.}

\subsubsection{Agent Lockdown}
\sectionstub{What a lockdown covers, including work the locked agent had
already started, and why it outranks any allowing rule.}
% MATERIAL: §3.5.38, §3.5.16.

\subsubsection{Kill Switch}
\sectionstub{Stopping a run, and why the design reports a confirmed stop
rather than a signal that was sent.}
% MATERIAL: §3.5.10, §3.5.21.
% DIVERGENCE from Chapter 1 (report/DIVERGENCES.md §1.4): §1.6 promises
%   a sub-second kill switch. The signal is milliseconds; the CONFIRMED
%   stop measured 2.2 s and 2.8 s on the laptop. Give both numbers and
%   say the design measures something stricter than Chapter 1 did.

\subsection{Human-in-the-Loop Approvals}
\label{sec:gov-approvals}
\sectionstub{What happens when the policy neither allows nor denies an
action and a person has to decide.}

\subsubsection{Escalation Routing}
\sectionstub{Who is entitled to answer a request, where it appears, and
what happens when nobody answers in time.}
% MATERIAL: §3.5.83.
% DIVERGENCE from Chapter 1 (report/DIVERGENCES.md §1.5): §1.6 puts the
%   timeout with Root and adds a per-user toggle. What was built sets
%   escalation, posture and timeout per AGENT, by an Administrator, and
%   has no per-user toggle.

\subsubsection{Persistent Approvals}
\sectionstub{What the ``always allow'' answer promises, what it does in
practice, and the rule request it files for an administrator.
Figure~\ref{fig:gov-alwaysallow}.}
% FIGURE F23 -> \ref{fig:gov-alwaysallow} (prints as Figure 3.11)
% MATERIAL: §3.5.81, §3.5.94.

\subsection{Runtime Management}
\label{sec:gov-runs}
\sectionstub{How running work is represented, bounded and governed,
including when the agent executes outside this process.}

\subsubsection{Task and Slot Model}
\sectionstub{The row an operator sees and the concurrency slot behind
it, and why the two are released at different moments.
Figure~\ref{fig:gov-taskslot}.}
% FIGURE F24 -> \ref{fig:gov-taskslot} (prints as Figure 3.12)
% MATERIAL: §3.5.82.

\subsubsection{Secondary Runtime Governance}
\sectionstub{The in-process path and the native harness, and the
two-layer permission that governs the second.
Figures~\ref{fig:gov-twopaths} and~\ref{fig:gov-codexgates}.}
% FIGURE F8  -> \ref{fig:gov-twopaths}   (prints as Figure 3.6)
% FIGURE F21 -> \ref{fig:gov-codexgates} (prints as Figure 3.9)
% MATERIAL: §3.5.15, §3.5.62.

\subsection{Tenancy and Agent Registry}
\label{sec:gov-tenant}
\sectionstub{One organization per installation, the storage that belongs
to it, and the registry that gives the layer a noun for an agent.
Figure~\ref{fig:gov-tenant}.}
% FIGURE F19 -> \ref{fig:gov-tenant} (prints as Figure 3.8)
% MATERIAL: §3.5.31, §3.5.33, §3.5.47, §3.5.89.

\subsection{Agent Lifecycle}
\label{sec:gov-lifecycle}
\sectionstub{Registration as the precondition for being allowed
anything; ownership and re-owning; and the two deletions, with what each
one leaves behind on the host.}
% MATERIAL: §3.5.33, §3.5.90, §3.5.91, §3.5.94.

\subsection{Management Interface}
\label{sec:gov-dashboard}
\sectionstub{The HTTP control plane and the dashboard behind it, what
its sections are for, and the rule that a control must never promise
what the gate cannot keep.}
% MATERIAL: §3.5.13, §3.5.37, §3.5.46, §3.5.80.

\subsection{System Security}
\label{sec:gov-security}
\sectionstub{The threats this design answers, and the three it does not:
the check-then-open window, a denied file still reachable on the native
harness, and the fact that the layer governs what an agent does rather
than why it does it. Figure~\ref{fig:gov-toctou}.}
% FIGURE F11 -> \ref{fig:gov-toctou} (prints as Figure 3.7)
% MATERIAL: §3.5.7, §3.5.61; mg/HANDOFF.md §7.
% CAUTION: do not write that the layer prevents prompt injection. It
%   contains the damage when persuasion succeeds. Chapter 2 sets this
%   up, so this section should answer Chapter 2 directly.

% =====================================================================
\section{Summary}
\label{sec:gov-designsummary}

\sectionstub{What the chapter established, the one requirement left
partially met, and the handover to Chapter 4, which tests the design
against Section 1.3 as its benchmark.}
